BFV supports exact arithmetic on plaintext polynomials modulo an integer plaintext modulus. Unlike approximate CKKS arithmetic, its decoded result is exact when every operation remains within the BFV correctness bound; encoding integers alone does not guarantee successful decryption for arbitrary noise or circuit depth. Applications must also choose a plaintext modulus and encoding that preserve their intended integer semantics. The reference scheme is that of Fan and Vercauteren~\cite{cryptoeprint:2012/144}.

In BFV, each plaintext is encrypted as an RLWE ciphertext. Therefore, BFV's ciphertext-to-ciphertext addition, ciphertext-to-plaintext addition, and ciphertext-to-plaintext multiplication are implemented based on GLWE's homomorphic addition and multiplication (as we learned in \autoref{part:generic-fhe}), with $k = 1$ to make GLWE an RLWE.




\begin{tcolorbox}[
    title = \textbf{Required Background},    % box title
    colback = white,    % light background; tweak to taste
    colframe = black,  % frame colour
    boxrule = 0.8pt,     % line thickness
    left = 1mm, right = 1mm, top = 1mm, bottom = 1mm % inner padding
]

\begin{itemize}
\item \autoref{sec:modulo}: \nameref{sec:modulo}
\item \autoref{sec:group}: \nameref{sec:group}
\item \autoref{sec:field}: \nameref{sec:field}
\item \autoref{sec:order}: \nameref{sec:order}
\item \autoref{sec:polynomial-ring}: \nameref{sec:polynomial-ring}
\item \autoref{sec:decomp}: \nameref{sec:decomp}
\item \autoref{sec:roots}: \nameref{sec:roots}
\item \autoref{sec:cyclotomic}: \nameref{sec:cyclotomic}
\item \autoref{sec:cyclotomic-polynomial-integer-ring}: \nameref{sec:cyclotomic-polynomial-integer-ring}
\item \autoref{sec:matrix}: \nameref{sec:matrix}
\item \autoref{sec:euler}: \nameref{sec:euler}
\item \autoref{sec:modulus-rescaling}: \nameref{sec:modulus-rescaling}
\item \autoref{sec:chinese-remainder}: \nameref{sec:chinese-remainder}
\item \autoref{sec:polynomial-interpolation}: \nameref{sec:polynomial-interpolation}
\item \autoref{sec:ntt}: \nameref{sec:ntt}
\item \autoref{sec:lattice}: \nameref{sec:lattice}
\item \autoref{sec:rlwe}: \nameref{sec:rlwe}
\item \autoref{sec:glwe}: \nameref{sec:glwe}
\item \autoref{sec:glwe-add-cipher}: \nameref{sec:glwe-add-cipher}
\item \autoref{sec:glwe-add-plain}: \nameref{sec:glwe-add-plain}
\item \autoref{sec:glwe-mult-plain}: \nameref{sec:glwe-mult-plain}
\item \autoref{subsec:modulus-switch-rlwe}: \nameref{subsec:modulus-switch-rlwe}
\item \autoref{sec:glwe-key-switching}: \nameref{sec:glwe-key-switching}
\end{itemize}
\end{tcolorbox}

\clearpage

\subsection{Single Value Encoding}
\label{subsec:bfv-single-encoding}

This binary-polynomial encoding requires $0\leq m<2^n$. Evaluation at $X=2$ preserves integer addition and multiplication only while the chosen integer coefficient lifts do not wrap modulo $t$ and the product degree is strictly less than $n$. It is not a ring homomorphism from $\mathbb Z_t[X]/(X^n+1)$ to the unbounded integers. For example, with $n=2$, the encoding $X$ of $2$ squares to $-1$ in the quotient ring, not the integer $4$.

BFV supports two encoding schemes: single value encoding and batch encoding. In this subsection, we will explain the single value encoding scheme. 



\begin{tcolorbox}[title={\textbf{\tboxlabel{\ref*{subsec:bfv-single-encoding}} BFV Encoding}}]

\textbf{\underline{Input Integer}:} Decompose an integer $m$ satisfying $0\leq m<2^n$ as follows:

$m = b_{n-1}\cdot 2^{n-1} + b_{n-2}\cdot 2^{n-2} + \cdots + b_1\cdot 2^1 + b_0\cdot 2^0$  \text { } , where each $b_i \in \{0, 1\}$

$ $


\textbf{\underline{Encoded Polynomial}:} $M(X) = b_0 + b_1X + b_2X^2 + \cdots + b_{n-1}X^{n-1} \in \mathcal{R}_{\langle n, t\rangle}$

$ $

\textbf{\underline{Decoding}:} $M(X=2) = m$

\end{tcolorbox}


Let's analyze whether the encoding scheme in Summary~\ref*{subsec:bfv-single-encoding} ensures correct decoding after addition and multiplication. This is equivalent to showing that:

$\text{Decode}(\sigma(m_1) + \sigma(m_2)) = m_1 + m_2$ \textcolor{red}{$\rhd$ where $\sigma$ is the encoding function}

$\text{Decode}(\sigma(m_1) \cdot \sigma(m_2)) = m_1 \cdot m_2$

Let $\sigma(m_1) = M_1(X)$ and $\sigma(m_2) = M_2(X)$. Then, 

$\text{Decode}(\sigma(m_1) + \sigma(m_2)) = \text{Decode}(M_1(X) + M_2(X))$


$= \text{Decode}(M_{1+2}(X))$ \textcolor{red}{$\rhd$ where $M_{1+2}(X) = M_1(X) + M_2(X)$}

$= M_{1+2}(2)$ \textcolor{red}{$\rhd$ since decoding is evaluating the polynomial at $X = 2$}

$= M_{1}(2) + M_{2}(2)$ \textcolor{red}{$\rhd$ since evaluating $M_{1+2}(X)$ at $X=2$ is computationally the same as splitting $M_{1+2}(X)$ into $M_{1}(X)$ and $M_{2}(X)$, evaluating $M_1(2)$ and $M_2(2)$ and summing them}

$ $

Similarly, the decoding preserves correctness over multiplication as well:

$\text{Decode}(\sigma(m_1) \cdot \sigma(m_2)) = \text{Decode}(M_1(X) \cdot M_2(X))$

$= \text{Decode}(M_{1\cdot2}(X))$ \textcolor{red}{$\rhd$ where $M_{1\cdot2}(X) = M_1(X) \cdot M_2(X)$}

$= M_{1\cdot2}(2)$ \textcolor{red}{$\rhd$ since decoding is evaluating the polynomial at $X = 2$}

$= M_{1}(2) \cdot M_{2}(2)$ \textcolor{red}{$\rhd$ since evaluating $M_{1\cdot2}(X)$ at $X=2$ is computationally the same as splitting $M_{1\cdot2}(X)$ into $M_{1}(X)$ and $M_{2}(X)$, evaluating $M_1(2)$ and $M_2(2)$ and multiplying them}




$ $

Therefore these integer identities hold under the stated no-wrap and degree conditions; they are not unconditional identities in the quotient ring.

$ $

The encoding of Summary~\ref*{subsec:bfv-single-encoding} preserves integer semantics only if every resulting polynomial has degree strictly below $n$ and its chosen integer coefficients do not wrap modulo $t$. At degree $n$, reduction by $X^n+1$ already changes evaluation at $2$.
Due to these constraints, the single value encoding is not a good choice for fully homomorphic encryption. Also, the single value encoding is computationally inefficient, because each polynomial can encode only a single value even if it holds $n$ coefficients. 


\subsection{Batch Encoding}
\label{subsec:bfv-batch-encoding}

For the scalar-slot formulas below, assume $n\geq2$ is a power of two and $t$ is prime with $2n\mid(t-1)$. Then $\omega=g^{(t-1)/(2n)}$ is a primitive $2n$-th root and the Vandermonde matrix is invertible. A prime-power modulus $t=p^r$ also supports these scalar slots when $p\equiv1\pmod{2n}$ and the roots are lifted to $\mathbb Z_{p^r}$; the exponent formula using $t-1$ is not valid in general for prime powers. Other parameters may give fewer slots over extension rings. In this chapter $W^T$ is the row-wise evaluation Vandermonde matrix, so $W$ is its transpose; this defines the matrix orientation explicitly.

While the single-value encoding scheme (\autoref{subsec:bfv-single-encoding}) encodes \& decodes each individual value one at a time, the batch encoding scheme does the same for a huge list of values simultaneously using a large dimensional vector. Therefore, batch encoding is more efficient than single-value encoding. Furthermore, batch-encoded values can be homomorphically added or multiplied simultaneously element-wise by vector-to-vector addition and Hadamard product. Therefore, the homomorphic operation of batch-encoded values can be processed more efficiently in a SIMD (single-instruction-multiple-data) manner than single-value encoded ones. 

BFV's encoding converts an $n$-slot integer input vector $\vec{v} = (v_0, v_1, v_2, \cdots v_{n-1})$ modulo $t$ into another $n$-slot vector $\vec{m} = (m_0, m_1, m_2, \cdots m_{n-1})$ modulo $t$, which are the coefficients of the encoded $(n-1)$-degree (or lesser-degree) polynomial $M(X) \in \mathbb{Z}_t[X] / (X^n + 1)$. 

$ $

\subsubsection{\textsf{Encoding\textsubscript{1}}}  
\label{subsubsec:bfv-encoding-1} 

In \autoref{subsec:poly-vector-transformation}, we learned that an $(n-1)$-degree (or lesser degree) polynomial can be isomorphically mapped to an $n$-slot vector based on the mapping $\sigma$ (we notate $\sigma_c$ in \autoref{subsec:poly-vector-transformation} as $\sigma$ for simplicity): 

$\sigma: M(X) \in \mathbb{Z}_t[X]/(X^n + 1) \longrightarrow  (M(\omega),M(\omega^3),M(\omega^5), \cdots, M(\omega^{2n-1})) \in \mathbb{Z}_t^n$

$ $

, which evaluates the polynomial $M(X)$ at $n$ distinct $(\mu=2n)$-th primitive roots of unity: $\omega, \omega^3, \omega^5, \cdots, \omega^{2n-1}$. Let $\vec{m}$ be a vector that contains $n$ coefficients of the polynomial $M(X)$. Then, we can express the mapping $\sigma$ as follows: 

$\vec{v} = W^T \cdot \vec{m}$

$ $

, where $W^T$ is as follows: 

$W^T = \begin{bmatrix}
1 & (\omega) & (\omega)^2 & \cdots & (\omega)^{n-1}\\
1 & (\omega^3) & (\omega^3)^2 & \cdots & (\omega^3)^{n-1}\\
1 & (\omega^5) & (\omega^5)^2 & \cdots & (\omega^5)^{n-1}\\
\vdots & \vdots & \vdots & \ddots & \vdots \\
1 & (\omega^{2n-1}) & (\omega^{2n-1})^2 & \cdots & (\omega^{2n-1})^{n-1}\\
\end{bmatrix}$  \textcolor{red}{ \text{ } \# $W^T$ is a transpose of $W$ described in \autoref{subsec:poly-vector-transformation}}

$ $

Note that the dot product between each row of $W^T$ and $\vec{m}$ computes the evaluation of $M(X)$ at each $X = \{w, w^3, w^5, \cdots, w^{2n-1}\}$. In the BFV encoding scheme, the \textsf{Encoding\textsubscript{1}} process encodes an $n$-slot input vector $\vec{v}$ $\in \mathbb{Z}_t^n$ into a plaintext polynomial $M(X) \in \mathbb{Z}_t[X] / (X^n + 1)$, and the \textsf{Decoding\textsubscript{2}} process decodes $M(X)$ back to $\vec{v}$. Since $W^T \cdot \vec{m}$ gives us $\vec{v}$ which is a decoding of $M(X)$, we call $W^T$ a decoding matrix. Meanwhile, the goal of \textsf{Encoding\textsubscript{1}} is to encode $\vec{v}$ into $M(X)$ so that we can do homomorphic computations based on $M(X)$. Given the relation $\vec{v} = W^T \cdot \vec{m}$, the encoding formula can be derived as follows:

$(W^T)^{-1} \cdot \vec{v} = (W^T)^{-1} W^T \cdot \vec{m}$

$\vec{m} = (W^T)^{-1} \cdot \vec{v}$

$ $

Therefore, we need to find out what $(W^T)^{-1}$ is, the inverse of $W^T$ as the encoding matrix. But we already learned from Theorem~\ref*{subsec:vandermonde-euler} (in \autoref{subsec:vandermonde-euler}) that $V^{-1} = \dfrac{V^T \cdot I_n^R}{n}$, where $V = W^T$ and $V^T = W$. In other words, $(W^T)^{-1} = \dfrac{W \cdot I_n^R}{n}$. Therefore, we can express the BFV encoding formula as: 

$\vec{m} = (W^T)^{-1} \cdot \vec{v} = \dfrac{W \cdot I_n^R \cdot \vec{v}}{n}$, \text{ } where

$ $

$W =  \begin{bmatrix}
1 & 1 & 1 & \cdots & 1\\
(\omega) & (\omega^3) & (\omega^5) & \cdots & (\omega^{2n-1})\\
(\omega)^2 & (\omega^3)^2 & (\omega^5)^2 & \cdots & (\omega^{2n-1})^2\\
\vdots & \vdots & \vdots & \ddots & \vdots \\
(\omega)^{n-1} & (\omega^3)^{n-1} & (\omega^5)^{n-1} & \cdots & (\omega^{2n-1})^{n-1}\\
\end{bmatrix}$

$= \begin{bmatrix}
1 & 1 & \cdots & 1 & 1 & \cdots & 1 & 1\\
(\omega) & (\omega^3) & \cdots & (\omega^{n - 1}) & (\omega^{-(n - 1)}) & \cdots & (\omega^{-3}) & (\omega^{-1})\\
(\omega)^2 & (\omega^3)^2 & \cdots & (\omega^{n - 1})^2 & (\omega^{-(n - 1)})^2 & \cdots & (\omega^{-3})^2 & (\omega^{-1})^2\\
\vdots & \vdots & \vdots & \ddots & \vdots \\
(\omega)^{n-1} & (\omega^3)^{n-1} & \cdots & (\omega^{n - 1})^{n-1} & (\omega^{-(n - 1)})^{n-1} & \cdots & (\omega^{-3})^{n-1} & (\omega^{-1})^{n-1}\\
\end{bmatrix}$

$ $

, where $\omega$ is a primitive $2n$-th root of unity modulo $t$ (which implies $t \equiv 1 \bmod 2n$). This implies that $\omega = g^{\frac{t - 1}{2n}} \bmod t$ ($g$ is a generator of $\mathbb{Z}_t^{\times}$ (see \autoref{subsubsec:poly-vector-transformation-modulus}).

In \autoref{subsec:poly-vector-transformation}, we learned that $W$ is a valid basis of the $n$-dimensional vector space. Therefore, $\dfrac{W \cdot I_n^R\cdot\vec{v}}{n} = \vec{m}$ is guaranteed to be a unique vector corresponding to each $\vec{v}$ in the $n$-dimensional vector space $\mathbb{Z}_t^{n}$ (refer to Theorem~\ref*{subsec:projection} in \autoref{subsec:projection}), and thereby the polynomial $M(X)$ comprising the $n$ elements of $\vec{m}$ as coefficients is a unique polynomial bi-jective to $\vec{v}$.

Note that by computing $\dfrac{W \cdot I_n^R \cdot \vec{v}}{n}$
, this is a linear bijection. The ring isomorphism identifies slotwise multiplication of $\vec v$ with negacyclic polynomial multiplication of $M$, not with componentwise multiplication of the coefficient vector $\vec m$.

\subsubsection{\textsf{Encoding\textsubscript{2}}}  

Once we have the $n$-slot vector $\vec{m}$, we scale (i.e., multiply) it by some scaling factor $\Delta = \left\lfloor \dfrac{q}{t} \right\rfloor$, where $q$ is the ciphertext modulus. We scale $\vec{m}$ by $\Delta$ and make it $\Delta \vec{m}$. 
The $n$ integers in $\Delta\vec{m}$ will be used as $n$ coefficients of the plaintext polynomial for RLWE encryption. The finally encoded plaintext polynomial $\Delta M = \sum\limits_{i=0}^{n-1}\Delta m_iX^i$. 


\subsubsection{\textsf{Decoding\textsubscript{1}}}  
\label{subsubsec:bfv-enc-dec-decoding1} 

The divisions by $\Delta$ in the encoding/decoding diagrams refer to the ideal, unreduced, noiseless encoded polynomial. For an actual ciphertext, first apply the scale-and-round decoder of \autoref{subsec:bfv-enc-dec}; only then apply the slot evaluation matrix to the recovered polynomial in $\mathcal R_{\langle n,t\rangle}$.

Once an RLWE ciphertext is (first-half) decrypted to $\Delta M = \sum\limits_{i=0}^{n-1}\Delta m_iX^i$, we compute $\dfrac{\Delta\vec{m}}{\Delta} = \vec{m}$.




\subsubsection{\textsf{Decoding\textsubscript{2}}}  


In \autoref{subsubsec:bfv-encoding-1}, we already derived the decoding formula that transforms an $(n-1)$-degree polynomial having integer modulo $t$ coefficients into an $n$-slot input vector as follows: 

$\vec{v} = W^T \cdot \vec{m}$


\subsubsection{Summary}
\label{subsubsec:bfv-encoding-summary}

\begin{tcolorbox}[title={\textbf{\tboxlabel{\ref*{subsubsec:bfv-encoding-summary}} BFV's Encoding and Decoding}}]



\textbf{\underline{Input}:} An $n$-slot integer modulo $t$ vector $\vec{v} = (v_0, v_1, \cdots, v_{n-1}) \in \mathbb{Z}_t^n$

\par\noindent\rule{\textwidth}{0.4pt}

\textbf{\underline{Encoding}} 
\begin{enumerate}

\item Convert $\vec{v} \in \mathbb{Z}_t^n$ into $\vec{m} \in \mathbb{Z}_t^n$ by applying the transformation $\vec{m} = n^{-1}\cdot W \cdot I_n^R \cdot \vec{v}$

, where $W$ is a basis of the $n$-dimensional vector space crafted as follows: 

$W =  \begin{bmatrix}
1 & 1 & 1 & \cdots & 1\\
(\omega) & (\omega^3) & (\omega^5) & \cdots & (\omega^{2n-1})\\
(\omega)^2 & (\omega^3)^2 & (\omega^5)^2 & \cdots & (\omega^{2n-1})^2\\
\vdots & \vdots & \vdots & \ddots & \vdots \\
(\omega)^{n-1} & (\omega^3)^{n-1} & (\omega^5)^{n-1} & \cdots & (\omega^{2n-1})^{n-1}\\
\end{bmatrix}$

$= \begin{bmatrix}
1 & 1 & \cdots & 1 & 1 & \cdots & 1 & 1\\
(\omega) & (\omega^3) & \cdots & (\omega^{n - 1}) & (\omega^{-(n - 1)}) & \cdots & (\omega^{-3}) & (\omega^{-1})\\
(\omega)^2 & (\omega^3)^2 & \cdots & (\omega^{n - 1})^2 & (\omega^{-(n - 1)})^2 & \cdots & (\omega^{-3})^2 & (\omega^{-1})^2\\
\vdots & \vdots & \vdots & \ddots & \vdots \\
(\omega)^{n-1} & (\omega^3)^{n-1} & \cdots & (\omega^{n - 1})^{n-1} & (\omega^{-(n - 1)})^{n-1} & \cdots & (\omega^{-3})^{n-1} & (\omega^{-1})^{n-1}\\
\end{bmatrix}$

$ $

, where $\omega$ is a primitive $2n$-th root of unity modulo $t$. This implies that $\omega = g^{\frac{t - 1}{2n}} \bmod t$ ($g$ is a generator of $\mathbb{Z}_t^{\times}$ (see \autoref{subsubsec:poly-vector-transformation-modulus}).

$ $

\item Convert $\vec{m}$ into a scaled integer vector $\Delta\vec{m}$, where $\Delta=\lfloor \dfrac{q}{t}\rfloor$ is a scaling factor. If $\Delta$ is too close to 1 (i.e., $t$ is too big), the \textbf{noise budget} will become too small (making decryption fail easily). If $\Delta$ is too close to $q$ (i.e., $t$ is too small), the \textbf{message capacity} will become too small (i.e., the plaintext modulus $t$ limits the range of values that can be encoded). The finally encoded plaintext polynomial $\Delta M = \sum\limits_{i=0}^{n-1} \Delta m_i X^i \text{ } \in \mathbb{Z}_q[X] / (X^n + 1)$. %The rounding of $\lceil \Delta \vec{m} \rfloor$ during the encoding process makes CKKS an approximate FHE scheme.

\end{enumerate}

\par\noindent\rule{\textwidth}{0.4pt}

\textbf{\underline{Decoding}:} From the plaintext polynomial $\Delta M = \sum\limits_{i=0}^{n-1}\Delta m_iX^i$, recover $\vec{m} = \dfrac{\Delta \vec{m}}{\Delta}$. Then,
compute $\vec{v} = W^T \cdot \vec{m}$.

\end{tcolorbox}


%\para{BFV's Encoding Range:} \autoref{tab:polynomial-encoding-capacity-integer} depicts the range of input values that can be encoded by the polynomials $M(X) \in \mathcal{R}_{\langle n, t \rangle}$ over $X \in \mathbb{Z}_t$. As illustrated in \autoref{tab:polynomial-encoding-capacity-integer}, the encoding range of integer inputs is determined by the size of the $n$ and $t$ parameters. Therefore, these two parameters should be chosen sufficiently large enough to encode all values that can be used by an application. 

However, Summary~\ref*{subsubsec:bfv-encoding-summary} is not the final version of BFV's batch encoding. In \autoref{subsec:bfv-rotation}, we will explain how to homomorphically rotate the input vector slots without decrypting the ciphertext that encapsulates it. To support such homomorphic rotation, we will need to slightly update the encoding scheme explained in Summary~\ref*{subsubsec:bfv-encoding-summary}. We will explain how to do this in \autoref{subsec:bfv-rotation}, and BFV's final encoding scheme is summarized in Summary~\ref*{subsubsec:bfv-rotation-summary} in \autoref{subsubsec:bfv-rotation-summary}.


\subsection{Encryption and Decryption}
\label{subsec:bfv-enc-dec}

We describe symmetric BFV encryption using the RLWE phase $B+AS$, the alternative sign convention of \autoref{subsec:glwe-alternative} with $k=1$ (\autoref{sec:rlwe}). Write $[U]_q$ for coefficientwise centered reduction and set $r_q=q-t\Delta$. The complete decoder $\textsf{Dec}$ uses $t/q$. Elsewhere in the phase calculations, $\textsf{RLWE}^{-1}$ denotes the intermediate noisy phase before this final scale-and-round step. Neither notation asserts fresh noise after a homomorphic operation.



\begin{tcolorbox}[title={\textbf{\tboxlabel{\ref*{subsec:bfv-enc-dec}} BFV Encryption and Decryption}}]

\textbf{\underline{Initial Setup}:} $\Delta=\left\lfloor\dfrac{q}{t}\right\rfloor \text{ is a plaintext scaling factor for polynomial encoding}, \text{ } S \xleftarrow{\$} \mathcal{R}_{\langle n, \textit{tern} \rangle}$

The basic plaintext modulus may be any integer $t\geq2$ with $q\gg t$. Batch encoding and the prime-power bootstrapping method impose the additional conditions stated in their respective sections. The coefficients of $S$ here are ternary, in $\{-1,0,1\}$.

\par\noindent\rule{\textwidth}{0.4pt}

\textbf{\underline{Encryption Input}:} $\Delta M \in \mathcal{R}_{\langle n, q \rangle}$, $A \xleftarrow{\$} \mathcal{R}_{\langle n, q \rangle}$, $E \xleftarrow{\chi_\sigma} \mathcal{R}_{\langle n, q \rangle}$


$ $

\begin{enumerate}


\item Compute $B = -A \cdot S + \Delta M + E \text{ } \in \mathcal{R}_{\langle n,q \rangle}$

\item $\textsf{RLWE}_{S,\sigma}(\Delta M + E) = (A, B) \text{ } \in \mathcal{R}_{\langle n,q \rangle}^2$ 

\end{enumerate}

\par\noindent\rule{\textwidth}{0.4pt}

\textbf{\underline{Decryption Input}:} $\textsf{ct} = (A, B) \text{ } \in \mathcal{R}_{\langle n,q \rangle}^2$ 

$\textsf{Dec}_S(A,B)=\left[\left\lceil\dfrac{t}{q}[B+AS]_q\right\rfloor\right]_t=M$.

(The noise $E = \sum\limits_{i=0}^{n-1}e_iX^i$ gets eliminated by the rounding process) 



$ $

\textbf{\underline{Conditions for Correct Decryption}:}

For a phase $\Delta M+E\pmod q$, a sufficient coefficientwise condition is $\|tE-r_qM\|_\infty<q/2$. In particular, $\|E\|_\infty<q/(2t)-r_q/2$ suffices when $M$ is centered with $\|M\|_\infty\leq t/2$, when the right side is positive. Plaintext reduction after homomorphic operations contributes additional remainder errors, as derived in \autoref{subsubsec:bfv-noise-analysis}. This is the BFV scale-and-round condition; the alternative $1/\Delta$ convention discussed in \autoref{subsubsec:lwe-noise-bound} must not be substituted without changing its assumptions.

\end{tcolorbox}


In this section, we will often write $\textsf{RLWE}_{S,\sigma}(\Delta  M + E)$ as $\textsf{RLWE}_{S,\sigma}(\Delta  M)$ for simplicity, because $\textsf{RLWE}_{S,\sigma}(\Delta M + E) \approx \textsf{RLWE}_{S,\sigma}(\Delta M)$ (i.e., they decrypt to the same message). Even in the case that we write $\textsf{RLWE}_{S,\sigma}(\Delta  M)$ instead of $\textsf{RLWE}_{S,\sigma}(\Delta  M + E)$, you should assume this as an encryption of $\Delta  M + E$ (i.e., the noise is included inside the scaled message). 


We will explain the conditions for BFV's correct decryption in more detail in \autoref{subsubsec:bfv-noise-analysis}.



\subsection{Ciphertext-to-Ciphertext Addition}
\label{subsec:bfv-add-cipher}

BFV's ciphertext-to-ciphertext addition uses RLWE's ciphertext-to-ciphertext addition scheme with the sign of the $A\cdot S$ term flipped in the encryption and decryption formula. Specifically, this is equivalent to the alternative GLWE version's (\autoref{subsec:glwe-alternative}) ciphertext-to-ciphertext addition scheme with $k = 1$. 


\begin{tcolorbox}[title={\textbf{\tboxlabel{\ref*{subsec:bfv-add-cipher}} BFV Ciphertext-to-Ciphertext Addition}}]
$\textsf{RLWE}_{S, \sigma}(\Delta M^{\langle 1 \rangle} + E^{\langle 1 \rangle} ) + \textsf{RLWE}_{S, \sigma}(\Delta M^{\langle 2 \rangle} + E^{\langle 2 \rangle}) $

$ = ( A^{\langle 1 \rangle}, \text{ } B^{\langle 1 \rangle}) + (A^{\langle 2 \rangle}, \text{ } B^{\langle 2 \rangle}) $

$ = ( A^{\langle 1 \rangle} + A^{\langle 2 \rangle}, \text{ } B^{\langle 1 \rangle} + B^{\langle 2 \rangle} ) $

$= \textsf{RLWE}_{S, \sigma}(\Delta(M^{\langle 1 \rangle} + M^{\langle 2 \rangle}) + E^{\langle 1 \rangle} + E^{\langle 2 \rangle})$
\end{tcolorbox}


\subsubsection{Noise Bound Analysis}
\label{subsubsec:bfv-noise-analysis}

Let $r_q=q-t\Delta$, and let the centered phase be $x=[\Delta M+E]_q=\Delta M+E-qK$. Multiplying by $t/q$ gives the exact identity
\[
\frac{t}{q}x=M+\frac{tE-r_qM}{q}-tK.
\]
The integral multiple $tK$ disappears after the final reduction modulo $t$. Therefore Summary~\ref*{subsec:bfv-enc-dec} in \autoref{subsec:bfv-enc-dec} decrypts correctly if
\[
\|tE-r_qM\|_\infty<q/2.
\]
This proof does not drop the inner reduction modulo $q$. Replacing $t/q$ by $1/\Delta$ does not give the same identity when $r_q\ne0$; it requires a separately specified decoder and representative convention.

For ciphertext addition, write $M^{\langle1\rangle}+M^{\langle2\rangle}=M^{\langle1+2\rangle}+tH$, where $M^{\langle1+2\rangle}$ is the chosen centered plaintext representative. Since $t\Delta=q-r_q$, the sum ciphertext has phase
\[
\Delta M^{\langle1+2\rangle}+E_{\rm add}\pmod q,
\qquad E_{\rm add}=E^{\langle1\rangle}+E^{\langle2\rangle}-r_qH.
\]
Its sufficient correctness condition is
\[
\bigl\|t(E^{\langle1\rangle}+E^{\langle2\rangle}-r_qH)-r_qM^{\langle1+2\rangle}\bigr\|_\infty<q/2.
\]
The term $-r_qH$ records plaintext wraparound and cannot be omitted. The same reasoning applies after multiplication or rotation once the effective message and complete error have been identified. It is consistent with the need for explicit remainder terms in \autoref{subsubsec:lwe-noise-bound}. A larger $q/t$ may increase the correctness budget, but security must be assessed for the complete parameter set.

\subsection{Ciphertext-to-Plaintext Addition}
\label{subsec:bfv-add-plain}

The added polynomial uses scale $\Delta$. If $M+\Lambda=M_{\rm out}+tH$, the effective error after plaintext reduction is $E-r_qH$, with $r_q=q-t\Delta$ as in \autoref{subsubsec:bfv-noise-analysis}.

BFV's ciphertext-to-plaintext addition uses RLWE's ciphertext-to-plaintext addition scheme with the sign of the $A\cdot S$ term flipped in the encryption and decryption formulas. Specifically, this is equivalent to the alternative GLWE version's (\autoref{subsec:glwe-alternative}) ciphertext-to-plaintext addition scheme (\autoref{sec:glwe-add-plain}) with $k = 1$.  


\begin{tcolorbox}[title={\textbf{\tboxlabel{\ref*{subsec:bfv-add-plain}} BFV Ciphertext-to-Plaintext Addition}}]
$\textsf{RLWE}_{S, \sigma}(\Delta M + E) + \Delta\Lambda $

$=  (A, \text{ } B) + \Delta\Lambda$

$=  (A, \text{ } B + \Delta\cdot\Lambda)$

$= \textsf{RLWE}_{S, \sigma}(\Delta (M + \Lambda) + E)$
\end{tcolorbox}



\subsection{Ciphertext-to-Plaintext Multiplication}
\label{subsec:bfv-mult-plain}

Here $\Lambda$ is an unscaled plaintext multiplier. If $M\Lambda=M_{\rm out}+tH$, the effective error is $\Lambda E-r_qH$. For polynomial multiplication, $\|\Lambda E\|_\infty\leq\|\Lambda\|_1\|E\|_\infty$; both this convolution bound and the remainder term enter the decoder condition.


BFV's ciphertext-to-plaintext multiplication uses RLWE's ciphertext-to-plaintext multiplication scheme with the sign of the $A\cdot S$ term flipped in the encryption and decryption formula. Specifically, this is equivalent to the alternative GLWE version's (\autoref{subsec:glwe-alternative}) ciphertext-to-plaintext multiplication scheme (\autoref{sec:glwe-mult-plain}) with $k = 1$.   


\begin{tcolorbox}[title={\textbf{\tboxlabel{\ref*{subsec:bfv-mult-plain}} BFV Ciphertext-to-Plaintext Multiplication}}]
$\textsf{RLWE}_{S, \sigma}(\Delta M + E) \cdot \Lambda$

$= (A, \text{ } B) \cdot \Lambda$

$= (A \cdot \Lambda, \text{ }  B \cdot \Lambda )$

$= \textsf{RLWE}_{S, \sigma}(\Delta (M \cdot \Lambda) + \Lambda E)$
\end{tcolorbox}




\subsection{Ciphertext-to-Ciphertext Multiplication}
\label{subsec:bfv-mult-cipher}

The four-stage presentation below is an explicitly bounded teaching variant: it uses an auxiliary modulus $Q=q\Delta$ and rescales by $1/\Delta$. The original FV multiplication scales integer-lifted tensor-product components by $t/q$ before relinearization; these factors coincide only when $t\mid q$. We keep this variant to explain modulus lifting, but retain its additional remainder and rounding errors. It is not a claim that every BFV implementation performs these four stages in this order.

\noindent \textbf{- Reference 1:} 
\href{https://www.inferati.com/blog/fhe-schemes-bfv}{Introduction to the BFV encryption scheme}~\cite{inferati-bfv}

\noindent \textbf{- Reference 2:} 
\href{https://eprint.iacr.org/2012/144.pdf}{Somewhat Practical Fully Homomorphic Encryption}~\cite{cryptoeprint:2012/144}

%\noindent \textbf{- Reference 3:} 
%\href{https://eprint.iacr.org/2016/510.pdf}{A Full RNS Variant $of FV like Somewhat Homomorphic Encryption Schemes}~\cite{cryptoeprint:2016/510}


Given two ciphertexts $\textsf{RLWE}_{S, \sigma}(\Delta M^{\langle 1\rangle })$ and $\textsf{RLWE}_{S, \sigma}(\Delta M^{\langle 2\rangle})$, the goal of ciphertext-to-ciphertext multiplication is to derive a new ciphertext whose decryption is $\Delta M^{\langle 1 \rangle} M^{\langle 2 \rangle}$. Ciphertext-to-ciphertext multiplication is more complex than ciphertext-to-plaintext multiplication. It comprises four steps: (1) \textsf{ModRaise}; (2) polynomial multiplication; (3) relinearization; and (4) rescaling. 

For better understanding, we will explain BFV's ciphertext-to-ciphertext multiplication based on the alternate version of RLWE (Theorem~\ref*{subsec:glwe-alternative} in \autoref{subsec:glwe-alternative}), where the sign of the $AS$ term is flipped in the encryption and decryption formulas.

\subsubsection{\textsf{ModRaise}}
\label{subsubsec:bfv-mult-cipher-modraise}

Choose integer coefficient representatives of $(A,B)$ and write their unreduced phase in $\mathbb Z[X]/(X^n+1)$ as
\[
AS+B=\Delta M+E+qK.
\]
Embedding these same coefficient representatives into modulus $Q=q\Delta$ gives phase $\Delta M+E+qK\pmod Q$, not an encryption of the unchanged $\Delta M$ with only error $E$. There is no general inequality $(q-1)n+(q-1)<Q$: it would require an additional relation between $n$ and $\Delta$. Nor is a no-wrap condition needed for the congruence. All products and phases in the next two stages are interpreted modulo $Q$; integer lifts and their multiples of $Q$ are accounted for explicitly when rescaling.

The first step of BFV's ciphertext-to-ciphertext multiplication is to \textit{mod-raise} the two input ciphertexts $(A^{\langle 1 \rangle}, B^{\langle 1 \rangle}) \bmod q$ and $(A^{\langle 2 \rangle}, B^{\langle 2 \rangle}) \bmod q$  
from $q \rightarrow Q$ (where $Q = q\cdot \Delta$) as follows: 

$(A^{\langle 1 \rangle}, B^{\langle 1 \rangle}) \bmod Q$

$(A^{\langle 2 \rangle}, B^{\langle 2 \rangle}) \bmod Q$

$ $

After \textsf{ModRaise}, the decryption of these two ciphertexts would be the following: 

$A^{\langle 1 \rangle} \cdot S + B^{\langle 1 \rangle} = \Delta M^{\langle 1 \rangle} + E^{\langle 1 \rangle} + K_1q \pmod Q$

$A^{\langle 2 \rangle} \cdot S + B^{\langle 2 \rangle} = \Delta M^{\langle 2 \rangle} + E^{\langle 2 \rangle} + K_2q \pmod Q$


$ $

Therefore, the \textit{mod-raised} ciphertexts have the following form: 

$\textsf{RLWE}_{S, \sigma}(\Delta M^{\langle 1 \rangle} + K_1q) = (A^{\langle 1 \rangle}, B^{\langle 1 \rangle}) \bmod Q$

$\textsf{RLWE}_{S, \sigma}(\Delta M^{\langle 2 \rangle} + K_2q) = (A^{\langle 2 \rangle}, B^{\langle 2 \rangle}) \bmod Q$


\subsubsection{Polynomial Multiplication}
\label{subsubsec:bfv-mult-cipher-multiplication}

Our next goal is to derive a new ciphertext which encrypts 
$(\Delta M^{\langle 1 \rangle} + E^{\langle 1 \rangle} + K_1q) \cdot (\Delta M^{\langle 2 \rangle} + E^{\langle 2 \rangle} + K_2q)$. 

First, we can derive the following relation: 

$(\Delta M^{\langle 1 \rangle} + E^{\langle 1 \rangle} + K_1q) \cdot (\Delta M^{\langle 2 \rangle} + E^{\langle 2 \rangle} + K_2q)$

$ = (A^{\langle 1 \rangle} \cdot S + B^{\langle 1 \rangle}) \cdot (A^{\langle 2 \rangle} \cdot S + B^{\langle 2 \rangle})$

$ = \underbrace{B^{\langle 1 \rangle}B^{\langle 2 \rangle}}_{D_0}  + \underbrace{(B^{\langle 2 \rangle}A^{\langle 1 \rangle} + B^{\langle 1 \rangle}A^{\langle 2 \rangle})}_{D_1} \cdot S + \underbrace{(A^{\langle 1 \rangle} \cdot A^{\langle 2 \rangle})}_{D_2} \cdot S \cdot S $



$= D_0 + D_1\cdot S + D_2\cdot S^2$



$ $

Meanwhile, we also have the following relations:

$\textsf{RLWE}_{S, \sigma}^{-1}(\Delta M^{\langle 1 \rangle} + K_1q) = \Delta M^{\langle 1 \rangle} + E^{\langle 1 \rangle} + K_1q$

$\textsf{RLWE}_{S, \sigma}^{-1}(\Delta M^{\langle 2 \rangle} + K_2q) = \Delta M^{\langle 2 \rangle} + E^{\langle 2 \rangle} + K_2q$

$ $

Combining all these, we reach the following relation: 

$\textsf{RLWE}_{S, \sigma}^{-1}(\Delta M^{\langle 1 \rangle} + K_1q) \cdot \textsf{RLWE}_{S, \sigma}^{-1}(\Delta M^{\langle 2 \rangle} + K_2q) = D_0 + D_1\cdot S + D_2\cdot S^2$

$ $

Notice that $D_0, D_1,$ and $D_2$ are known values as ciphertext components, whereas $S$ is only known to the private key owner. Therefore, we can view $D_0 + D_1\cdot S + D_2\cdot S^2$ as a decryption formula such that given the ciphertext components $D_0, D_1, D_2$ and the private key $S$, one can derive $(\Delta M^{\langle 1 \rangle} + E^{\langle 1 \rangle} + K_1q) \cdot (\Delta M^{\langle 2 \rangle} + E^{\langle 2 \rangle} + K_2q)$. In other words, we can let $(D_0, D_1, D_2)$ be a new form of ciphertext which can be decrypted by $S$ into $(\Delta M^{\langle 1 \rangle} + E^{\langle 1 \rangle} + K_1q) \cdot (\Delta M^{\langle 2 \rangle} + E^{\langle 2 \rangle} + K_2q)$. 

The triplet $(D_0,D_1,D_2)$ is a degree-two ciphertext in the secret $S$, so decryption uses the square $S^2$, not a square root. Further multiplication adds secret degrees; repeated squaring doubles the degree at each step. Relinearization returns to two components and degree one, at the cost of evaluation-key noise.



\subsubsection{Relinearization}
\label{subsubsec:bfv-mult-cipher-relinearization}

For arbitrary $Q$, this subsection uses the integer gadget $g_i=\beta^{i-1}$ with $l=\lceil\log_\beta Q\rceil$ and exact coefficient decomposition $D_2\equiv\sum_{i=1}^lD_{2,i}g_i\pmod Q$. This is the increasing-weight convention, rather than the descending $Q/\beta^i$ convention; its matching RLev rows encrypt $g_iS^2$. Let $E_{\rm rl}=\sum_iD_{2,i}E_i'$ denote their accumulated error. The relinearized phase equals the original phase plus $E_{\rm rl}$, not exactly the original phase.

Relinearization turns the triplet $(D_0,D_1,D_2)\in\mathcal R_{\langle n,Q\rangle}^3$ into two components. With $\textsf{ct}_\alpha=(D_1,D_0)$ and the gadget construction of $\textsf{ct}_\beta$, the relation is
\[
\phi_S(\textsf{ct}_\alpha+\textsf{ct}_\beta)\equiv D_0+D_1S+D_2S^2+E_{\rm rl}\pmod Q.
\]

In $D_0+D_1S+D_2S^2$, the first two terms form the synthetic ciphertext $\textsf{ct}_\alpha=(D_1,D_0)$. The gadget construction produces $\textsf{ct}_\beta$ with phase $\phi_S(\textsf{ct}_\beta)\equiv D_2S^2+E_{\rm rl}\pmod Q$. Publishing encryptions of functions of the secret, such as $S^2$, also requires the evaluation-key security assumption of the chosen construction; the algebra below is not a proof of that assumption.

A naive way of creating a ciphertext that encrypts $D_2 \cdot S^2$ is as follows: we encrypt $S^2$ into an RLWE ciphertext as $\textsf{RLWE}_{S, \sigma}(S^2) = (A^{\langle s \rangle}, B^{\langle s \rangle})$ such that $A^{\langle s \rangle}\cdot S + B^{\langle s \rangle} = S^2 + E^{\langle s \rangle} \bmod Q$ (where the ciphertext modulus is $Q$ and the plaintext scaling factor $\Delta = 1$). Then, we perform a ciphertext-to-plaintext multiplication (\autoref{sec:glwe-mult-plain}) with $D_2$, treating $D_2$ as a plaintext polynomial in modulo $Q$. However, this approach does not work in practice, because computing $D_2 \cdot \textsf{RLWE}_{S, \sigma}(S^2)$ generates a huge noise as follows:

$D_2 \cdot (A^{\langle s \rangle}, B^{\langle s \rangle}) = (D_2\cdot A^{\langle s \rangle}, D_2 \cdot B^{\langle s \rangle})$

$ $

, whose decryption is:

$D_2\cdot A^{\langle s \rangle} \cdot S + D_2\cdot B^{\langle s \rangle} = D_2 \cdot S^2 + D_2\cdot E^{\langle s \rangle} \pmod Q$ 

$ $

. In the above decrypted expression $D_2S^2 + D_2E^{\langle s \rangle} \bmod Q$, the term $D_2S^2$ is okay to be reduced modulo $Q$, because this term is originally allowed to be reduced modulo $Q$ in the final decryption formula $D_0 + D_1S + D_2S^2 \bmod Q$ as well. However, the problematic term is the noise $D_2\cdot E^{\langle s \rangle}$, because its coefficients can be any value in $[0, Q - 1]$ (since each coefficient of polynomial $D_2 = A^{\langle 1 \rangle} A^{\langle 2 \rangle}$ can be any value in $[0, Q - 1]$). Such a huge noise is not allowed for correct final decryption. 

To avoid this noise issue, an improved solution is to express the RLWE ciphertext that encrypts $D_2 S^2$ as additions of multiple RLWE ciphertexts with small noises by using the gadget decomposition technique (\autoref{subsec:gadget-decomposition}). For this, we use an RLev ciphertext (\autoref{sec:glev}) that encrypts $S^2$. Suppose our gadget vector is $\vec{g} = \Bigg(\beta^{0}, \beta^{1}, \beta^{2}, \cdots, \beta^{l-1}\Bigg)$. Remember that our goal is to find $\textsf{ct}_\beta = \textsf{RLWE}_{S, \sigma}( S^2 \cdot D_2)$ given known $D_2$, unknown $S$, and known $\textsf{RLev}_{S, \sigma}^{\beta, l}(S^2) = \left\{\textsf{RLWE}_{S, \sigma}\left(\beta^{i-1}\cdot S^2\right)\right\}_{i=1}^{l}$. Then, we can derive $\textsf{ct}_\beta$ as follows:

$\textsf{ct}_\beta = \textsf{RLWE}_{S, \sigma}(S^2 \cdot D_2)$

$= \textsf{RLWE}_{S, \sigma}\left( S^2 \cdot \left(D_{2,1}\beta^{0} + D_{2,2}\beta^{1} + \cdots D_{2,l}\beta^{l-1}\right)\right)$ \textcolor{red}{  $\rhd$ by decomposing $D_2$}

$= \textsf{RLWE}_{S, \sigma}\left( S^2 \cdot D_{2,1} \cdot \beta^{0}\right) + \textsf{RLWE}_{S, \sigma}\left( S^2 \cdot D_{2,2} \cdot \beta^{1}\right) + \cdots + \textsf{RLWE}_{S, \sigma}\left( S^2 \cdot D_{2,l} \cdot \beta^{l-1}\right)$

$= D_{2,1}\cdot\textsf{RLWE}_{S, \sigma}\left( S^2 \cdot \beta^{0}\right) + D_{2,2}\cdot\textsf{RLWE}_{S, \sigma}\left( S^2 \cdot \beta^{1}\right) + \cdots + D_{2,l}\cdot\textsf{RLWE}_{S, \sigma}\left( S^2 \cdot \beta^{l-1}\right)$
\textcolor{red}{ $\rhd$ where each \textsf{RLWE} ciphertext is an encryption of $S^2\beta^{0}, S^2\beta^{1}, \cdots, S^2\beta^{l-1}$ as plaintext with the plaintext scaling factor $\Delta = 1$}

$ $

$= \bm{\langle} \textsf{Decomp}^{\beta, l}(D_2), \text{ } \textsf{RLev}_{S, \sigma}^{\beta, l}( S^2) \bm{\rangle}$ \textcolor{red}{  $\rhd$ inner product of \textsf{Decomp} and \textsf{RLev} treating them as vectors}

$ $

If we decrypt the above, we get the following:

$\textsf{RLWE}_{S, \sigma}^{-1}(\textsf{ct}_\beta = \bm{\langle} \textsf{Decomp}^{\beta, l}(D_2), \text{ } \textsf{RLev}_{S, \sigma}^{\beta, l}( S^2) \bm{\rangle} \bm{)}$ \textcolor{red}{ $\rhd$ the scaling factors of $\textsf{RLev}_{S, \sigma}^{\beta, l}( S^2)$ are all 1}

$= D_{2,1}\cdot\left(E_1' +  S^2\beta^{0}\right) + D_{2,2}\cdot\left(E_2' +  S^2\beta^{1}\right) + \cdots + D_{2,l}\cdot\left(E_l' +  S^2\beta^{l-1}\right)$ \textcolor{red}{ \text{ } \# where each $E_i'$ is a noise embedded in $\textsf{RLWE}_{S, \sigma}\left(S^2\cdot\beta^{i-1}\right)$}

$ $

$= \sum\limits_{i=1}^{l} (E_i'\cdot D_{2,i}) +  S^2\cdot\left(D_{2,1}\beta^{0} + D_{2,2}\beta^{1} + \cdots + D_{2, l}\beta^{l-1}\right)$

$= \sum\limits_{i=1}^{l} \epsilon_{i} + D_2\cdot S^2$ \textcolor{red}{ \text{ } \# where each $\epsilon_i = E_i'\cdot D_{2,i}$}

$\equiv D_2S^2+E_{\rm rl}\pmod Q$, with $\|E_{\rm rl}\|_\infty\leq\sum_{i=1}^l\|D_{2,i}\|_1\|E_i'\|_\infty$. Small digits give a useful bound, not a strict comparison with every possible single-key noise realization.

$ $

Therefore, we get the following comprehensive relation:


$\textsf{RLWE}_{S, \sigma}^{-1}(\Delta M^{\langle 1 \rangle} + K_1q) \cdot \textsf{RLWE}_{S, \sigma}^{-1}(\Delta M^{\langle 2 \rangle} + K_2q) \bmod Q$

$= (\Delta M^{\langle 1 \rangle} + E^{\langle 1 \rangle} + K_1q) \cdot (\Delta M^{\langle 2 \rangle} + E^{\langle 2 \rangle} + K_2q) \bmod Q$

$ = (A^{\langle 1 \rangle} \cdot S + B^{\langle 1 \rangle}) \cdot (A^{\langle 2 \rangle} \cdot S + B^{\langle 2 \rangle})  \bmod Q$

$ = D_0 + D_1\cdot S + D_2\cdot S^2  \bmod Q$ \textcolor{red}{ $\rhd$ $D_0 = B^{\langle 1 \rangle} B^{\langle 2 \rangle}$, \text{ } $D_1 = A^{\langle 1 \rangle} B^{\langle 2 \rangle} + A^{\langle 2 \rangle} B^{\langle 1 \rangle}$, \text{ } $D_2 = A^{\langle 1 \rangle} A^{\langle 2 \rangle}$}

$ = \textsf{RLWE}_{S, \sigma}^{-1}(\textsf{ct}_\alpha) + \textsf{RLWE}_{S, \sigma}^{-1}(\textsf{ct}_\beta) - \sum\limits_{i=1}^{l} (E_i'\cdot D_{2,i})  \bmod Q$ 

\textcolor{red}{ $\rhd$ $\textsf{ct}_\alpha = (D_1, D_0) = (A_\alpha, B_\alpha)$, \text{ } $\textsf{ct}_\beta = \langle \textsf{Decomp}^{\beta, l}(D_2), \textsf{RLev}_{S, \sigma}^{\beta, l}(S^2) \rangle = (A_\beta, B_\beta)$}

$ $

$ = \textsf{RLWE}_{S, \sigma}^{-1}(\textsf{ct}_\alpha + \textsf{ct}_\beta) - \sum\limits_{i=1}^{l} (E_i'\cdot D_{2,i})  \bmod Q$

$ = \textsf{RLWE}_{S, \sigma}^{-1}\bm((A_{\alpha+\beta}, B_{\alpha+\beta})\bm) - \sum\limits_{i=1}^{l} (E_i'\cdot D_{2,i}) \bmod Q$  \textcolor{red}{ $\rhd$ $A_{\alpha+\beta} = A_{\alpha} + A_\beta$, \text{ } $B_{\alpha + \beta} = B_\alpha + B_\beta$}


$ $

From the above, we extract the following relation: 

$(\Delta M^{\langle 1 \rangle} + E^{\langle 1 \rangle} + K_1q) \cdot (\Delta M^{\langle 2 \rangle} + E^{\langle 2 \rangle} + K_2q) \bmod Q$

$=  \Delta^2M^{\langle 1 \rangle}M^{\langle 2 \rangle} + \Delta\cdot (M^{\langle 1 \rangle}E^{\langle 2 \rangle} + M^{\langle 2 \rangle}E^{\langle 1 \rangle}) + q\cdot(\Delta M^{\langle 1 \rangle}K_2 + \Delta M^{\langle 2 \rangle}K_1 + E^{\langle 1 \rangle}K_2 + E^{\langle 2 \rangle}K_1) + K_1 K_2 q^2 + E^{\langle 1 \rangle} E^{\langle 2 \rangle} \bmod Q$

$= \textsf{RLWE}_{S, \sigma}^{-1}\bm((A_{\alpha+\beta}, B_{\alpha+\beta})\bm) - \sum\limits_{i=1}^{l} (E_i'\cdot D_{2,i}) \bmod Q$  

$ $

We can re-write the above relation as follows:


$\textsf{RLWE}_{S, \sigma}^{-1}\bm((A_{\alpha+\beta}, B_{\alpha+\beta})\bm) = A_{\alpha+\beta}\cdot S + B_{\alpha+\beta} \bmod Q$

$ = \Delta^2M^{\langle 1 \rangle}M^{\langle 2 \rangle} + \Delta\cdot (M^{\langle 1 \rangle}E^{\langle 2 \rangle} + M^{\langle 2 \rangle}E^{\langle 1 \rangle}) + q\cdot(\Delta M^{\langle 1 \rangle}K_2 + \Delta M^{\langle 2 \rangle}K_1 + E^{\langle 1 \rangle}K_2 + E^{\langle 2 \rangle}K_1) + K_1 K_2 q^2 + E^{\langle 1 \rangle} E^{\langle 2 \rangle} + \sum\limits_{i=1}^{l} (E_i'\cdot D_{2,i}) \bmod Q$ 

$ $

The expression contains the product message, the lift terms $qK_i$, the input-error products and the relinearization error. The next subsection follows all of these terms through the $1/\Delta$ rescaling of this teaching variant and applies the complete BFV decoder condition; approximate equality with a clean encryption is not assumed.


\subsubsection{Rescaling}
\label{subsubsec:bfv-mult-cipher-rescaling}

Let $(A_*,B_*)=(A_{\alpha+\beta},B_{\alpha+\beta})$ be the relinearized ciphertext modulo $Q=q\Delta$, and choose integer coefficient representatives. Define
\[
C_A=\Bigl\lceil A_*/\Delta\Bigr\rfloor,\quad C_B=\Bigl\lceil B_*/\Delta\Bigr\rfloor,
\quad R_A=C_A-A_*/\Delta,\quad R_B=C_B-B_*/\Delta.
\]
Each coefficient of $R_A,R_B$ has absolute value at most $1/2$. Write $L_i=\Delta M^{\langle i\rangle}+E^{\langle i\rangle}+qK_i$ for the integer input phases, and retain $E_{\rm rl}$ from \autoref{subsubsec:bfv-mult-cipher-relinearization}. Then
\[
A_*S+B_*=L_1L_2+E_{\rm rl}+QK_*,
\qquad C_AS+C_B\equiv\frac{L_1L_2+E_{\rm rl}}{\Delta}+R_AS+R_B\pmod q.
\]
The equality uses $Q/\Delta=q$, so the discarded multiple of $Q$ becomes an integral multiple of $q$. Set $r_q=q-t\Delta$. Expanding the product and removing integral multiples of $q$ gives
\[
C_AS+C_B\equiv\Delta M^{\langle1\rangle}M^{\langle2\rangle}+\eta\pmod q,
\]
\[
\begin{split}
\eta={}&M^{\langle1\rangle}E^{\langle2\rangle}+M^{\langle2\rangle}E^{\langle1\rangle}
+\frac q\Delta(E^{\langle1\rangle}K_2+E^{\langle2\rangle}K_1)\\
&+\frac{q r_q}{\Delta}K_1K_2
+\frac{E^{\langle1\rangle}E^{\langle2\rangle}+E_{\rm rl}}{\Delta}+R_AS+R_B.
\end{split}
\]
In particular, $q^2/\Delta=qt+q r_q/\Delta$. The second term is not generally negligible and must not be discarded as an approximately zero error. The expression for $\eta$ may use rational summands, but its complete value agrees with the integer output phase modulo $q$.

To decode, write $M^{\langle1\rangle}M^{\langle2\rangle}=M_{\rm out}+tH$ using the chosen centered plaintext representative. The effective output error is $E_{\rm out}=\eta-r_qH$. A sufficient correctness condition for the standard BFV decoder is
\[
\|t(\eta-r_qH)-r_qM_{\rm out}\|_\infty<q/2.
\]
When $t\mid q$, the $r_q$ terms vanish, but message-noise products, evaluation-key error and component rounding remain.

\para{Noise Analysis:} All products above are negacyclic polynomial products. Use $\|UV\|_\infty\leq\|U\|_1\|V\|_\infty\leq n\|U\|_\infty\|V\|_\infty$, not a coefficientwise scalar bound. Also $\|R_AS+R_B\|_\infty\leq(\|S\|_1+1)/2$. Bound each remaining summand of $\eta$ with these inequalities and apply the triangle inequality before comparing to the final threshold. No unconditional successful-multiplication claim follows merely from $q\gg t$ or $Q=q\Delta$.

\para{Modulus Switch v.s. Rescaling:} In this particular construction, scaling components by $1/\Delta=q/Q$ changes the auxiliary modulus from $Q$ back to $q$ and reduces the message scale. The terms describe different goals and are not universally interchangeable. Standard FV multiplication can restore the message scale while retaining the external ciphertext modulus, by performing scale-and-round on unreduced integer products.

\subsubsection{Summary}
\label{subsubsec:bfv-mult-cipher-summary}

To put all things together, BFV's ciphertext-to-ciphertext multiplication is summarized as follows:

\begin{tcolorbox}[title={\textbf{\tboxlabel{\ref*{subsubsec:bfv-mult-cipher-summary}} BFV's Ciphertext-to-Ciphertext Multiplication}}]


Suppose we have the following two RLWE ciphertexts:

$\textsf{RLWE}_{S, \sigma}(\Delta M^{\langle 1 \rangle}) = (A^{\langle 1 \rangle}, B^{\langle 1 \rangle}) \bmod q$, \text{ } where $B^{\langle 1 \rangle} = -A^{\langle 1 \rangle} \cdot S + \Delta M^{\langle 1 \rangle} + E^{\langle 1 \rangle} \bmod q$

$\textsf{RLWE}_{S, \sigma}(\Delta M^{\langle 2 \rangle}) = (A^{\langle 2 \rangle}, B^{\langle 2 \rangle})  \bmod q$, \text{ } where $B^{\langle 2 \rangle} = -A^{\langle 2 \rangle} \cdot S + \Delta M^{\langle 2 \rangle} + E^{\langle 2 \rangle} \bmod q$

$ $

Multiplication between these two ciphertexts is performed as follows:

$ $

\begin{enumerate}

\item \textbf{\underline{\textsf{ModRaise}}}

Forcibly raise the modulus of the ciphertexts $(A^{\langle 1 \rangle}, B^{\langle 1 \rangle}) \bmod q$ and $(A^{\langle 2 \rangle}, B^{\langle 2 \rangle}) \bmod q$ to $Q$ (where $Q = q \cdot \Delta$) as follows:

$(A^{\langle 1 \rangle}, B^{\langle 1 \rangle}) \bmod Q$

$(A^{\langle 2 \rangle}, B^{\langle 2 \rangle}) \bmod Q$


$ $

\item \textbf{\underline{Multiplication}}

Compute the following polynomial multiplications in modulo $Q$: 

$D_0 = B^{\langle 1 \rangle}B^{\langle 2 \rangle} \bmod Q$

$D_1 = B^{\langle 2 \rangle}A^{\langle 1 \rangle} +  B^{\langle 1 \rangle}A^{\langle 2 \rangle} \bmod Q$

$D_2 = A^{\langle 1 \rangle} \cdot A^{\langle 2 \rangle} \bmod Q$


$ $

\item \textbf{\underline{Relinearization}} 


Compute the following:

$ \textsf{ct}_\alpha = (D
_1, D_0)$

$ \textsf{ct}_\beta = \bm{\langle} \textsf{Decomp}^{\beta, l}(D_2), \text{ } \textsf{RLev}_{S, \sigma}^{\beta, l}( S^2) \bm{\rangle}$.

$ \textsf{ct}_{\alpha + \beta} = \textsf{ct}_{\alpha} + \textsf{ct}_{\beta}$

$ $

Then, the following property holds: 

$\phi_S(\textsf{ct}_{\alpha+\beta})\equiv(\Delta M^{\langle1\rangle}+E^{\langle1\rangle}+qK_1)(\Delta M^{\langle2\rangle}+E^{\langle2\rangle}+qK_2)+E_{\rm rl}\pmod Q$.

$ $




\item \textbf{\underline{Rescaling}}

Update $\textsf{ct}_{\alpha + \beta} = (A_{\alpha + \beta}, B_{\alpha + \beta}) \bmod Q$ to $\textsf{ct'}_{\alpha + \beta} = \left(\left\lceil\dfrac{A_{\alpha + \beta}}{\Delta}\right\rfloor, \left\lceil\dfrac{B_{\alpha + \beta}}{\Delta}\right\rfloor\right) \bmod q$. 

$ $

This plaintext rescaling process can be also viewed as a modulus switch of the ciphertext $\textsf{ct}_{\alpha + \beta}$ from $Q \rightarrow q$. 


$ $

\end{enumerate}

Note that after the ciphertext-to-ciphertext multiplication, the plaintext scaling factor $\Delta=\left\lfloor\dfrac{q}{t}\right\rfloor$, the ciphertext modulus $q$, and the private key $S$ stay the same as before.



\end{tcolorbox}


\para{The Purpose of \textsf{ModRaise}:} Reducing a tensor product modulo $q$ before scaling can destroy information needed for scale-and-round. This teaching variant instead computes modulo $Q=q\Delta$ because a discarded $Q$ multiple becomes a $q$ multiple after division by $\Delta$. It does not require every plaintext product to lie below $Q$; that claim would in general fail for polynomial convolution and large plaintexts. The lifting terms $qK_i$ and the complete error derived in \autoref{subsubsec:bfv-mult-cipher-rescaling} must still be included.

$ $

\para{Swapping the Order of \textsf{Relinearization} and \textsf{Rescaling}:} Both orders can be designed, but their evaluation keys, moduli and rounding-error bounds differ. They are not a literal swap using the same keys and error proof. Standard FV commonly scales the degree-two integer product before relinearizing at the target modulus.



\subsubsection{Application to TFHE's Ciphertext-to-Ciphertext Multiplication}
\label{subsubsec:bfv-mult-cipher-tfhe}

This subsection deliberately uses the alternative LWE phase $b+\langle\vec a,\vec s\rangle$. To apply it to the TFHE convention $b-\langle\vec a,\vec s\rangle$ of \autoref{subsec:tfhe-enc-dec}, first negate the public mask and negate it back on output. The tensor-product construction is a leveled LWE multiplication alternative, not the external product used by TFHE bootstrapping. It needs evaluation keys for the products $s_i s_j$, appropriate parameters and the full rescaling error bounds; it does not make arbitrary LWE multiplication noise-free.


In \autoref{subsec:tfhe-mult-cipher}, the external product takes an LWE ciphertext and a GSW ciphertext. Obtaining the latter from an arbitrary LWE ciphertext may require circuit bootstrapping. The tensor construction below instead requires a different, quadratic-size evaluation key and a rescaling budget. This establishes an alternative operation, without asserting a performance advantage over TFHE's external product or programmable bootstrapping. We impose $t\mid q$ and $\Delta=q/t$, as in the TFHE encoding of this book. Given two LWE ciphertexts, their integer lifts satisfy:

$\textsf{LWE}_{\vec{s}, \sigma}(\Delta m^{\langle 1 \rangle} + e^{\langle 1 \rangle}) = (\vec{a}^{\langle 1 \rangle}, b^{\langle 1 \rangle }) \Rightarrow \,\,\, \vec{a}^{\langle 1 \rangle} \cdot \vec{s} + b^{\langle 1 \rangle} = \Delta m^{\langle 1 \rangle} + e^{\langle 1 \rangle} + k_1q$

$\textsf{LWE}_{\vec{s}, \sigma}(\Delta m^{\langle 2 \rangle} + e^{\langle 2 \rangle})  = (\vec{a}^{\langle 2 \rangle}, b^{\langle 2 \rangle}) \Rightarrow \,\,\, \vec{a}^{\langle 2 \rangle} \cdot \vec{s} + b^{\langle 2 \rangle} = \Delta m^{\langle 2 \rangle} + e^{\langle 2 \rangle} + k_2q$

$ $

Multiplying these two equations yields the following relation: 


$(\Delta m^{\langle 1 \rangle} + e^{\langle 1 \rangle} + k_1q) \cdot (\Delta m^{\langle 2 \rangle} + e^{\langle 2 \rangle} + k_2q)$

$ = (\vec{a}^{\langle 1 \rangle} \cdot \vec{s} + b^{\langle 1 \rangle}) \cdot (\vec{a}^{\langle 2 \rangle} \cdot \vec{s} + b^{\langle 2 \rangle})$ 

$ = b^{\langle 1 \rangle}b^{\langle 2 \rangle}  + (b^{\langle 2 \rangle}\vec{a}^{\langle 1 \rangle} + b^{\langle 1 \rangle}\vec{a}^{\langle 2 \rangle}) \cdot \vec{s} + (\vec{a}^{\langle 1 \rangle} \cdot \vec{s}) \cdot (\vec{a}^{\langle 2 \rangle} \cdot \vec{s}) $



$ = \underbrace{b^{\langle 1 \rangle}b^{\langle 2 \rangle}}_{d_0}  + \underbrace{(b^{\langle 2 \rangle}\vec{a}^{\langle 1 \rangle} + b^{\langle 1 \rangle}\vec{a}^{\langle 2 \rangle})}_{d_1} \cdot \vec{s} + \underbrace{(\vec{a}^{\langle 1 \rangle} \otimes \vec{a}^{\langle 2 \rangle})}_{\vec{d}_2} \cdot (\vec{s} \otimes \vec{s}) $  


$= d_0 + d_1\cdot \vec{s} + \vec{d}_2\cdot \vec{s}^{\otimes}$ \textcolor{red}{ $\rhd$ where $\vec{s}^{\otimes} = \vec{s} \otimes \vec{s}$}

$ $

, where $\otimes$ denotes an outer product of two vectors. For example, given two $n$-length vectors $\vec{v}$ and $\vec{u}$, $\vec{v} \otimes \vec{u}$ is equivalent to an $n^2$-length vector that concatenates the following $n$ distinct $n$-length vectors: $v_0\cdot\vec{u}, \,\, v_1\cdot\vec{u}, \,\, \cdots v_{n-1}\cdot\vec{u}$. Notice that the above relation is similar to that we derived in BFV's ciphertext-to-ciphertext multiplication. Therefore, similar to the remaining steps in BFV, we can relinearize and rescale this LWE term as follows:


\begin{tcolorbox}[title={\textbf{\tboxlabel{\ref*{subsubsec:bfv-mult-cipher-tfhe}} Leveled LWE Tensor Multiplication}}]


Suppose we have the following two LWE ciphertexts:

$\textsf{LWE}_{\vec{s}, \sigma}(\Delta M^{\langle 1 \rangle}) = (\vec{a}^{\langle 1 \rangle}, b^{\langle 1 \rangle}) \bmod q$, \text{ } where $b^{\langle 1 \rangle} = -\vec{a}^{\langle 1 \rangle} \cdot \vec{s} + \Delta m^{\langle 1 \rangle} + e^{\langle 1 \rangle} \bmod q$

$\textsf{LWE}_{\vec{s}, \sigma}(\Delta M^{\langle 2 \rangle}) = (\vec{a}^{\langle 2 \rangle}, b^{\langle 2 \rangle}) \bmod q$, \text{ } where $b^{\langle 2 \rangle} = -\vec{a}^{\langle 2 \rangle} \cdot \vec{s} + \Delta m^{\langle 2 \rangle} + e^{\langle 2 \rangle} \bmod q$

$ $

Multiplication between these two LWE ciphertexts is performed as follows:

$ $

\begin{enumerate}

\item \textbf{\underline{\textsf{ModRaise}}}

Interpret the chosen integer component lifts of $(\vec{a}^{\langle 1 \rangle}, b^{\langle 1 \rangle})$ and $(\vec{a}^{\langle 2 \rangle}, b^{\langle 2 \rangle})$ modulo $Q=q\Delta$. Their phases still include the terms $qk_1$ and $qk_2$; this is not fresh encryption of the same scaled messages at the larger modulus:

$(\vec{a}^{\langle 1 \rangle}, b^{\langle 1 \rangle}) \bmod Q$

$(\vec{a}^{\langle 2 \rangle}, b^{\langle 2 \rangle}) \bmod Q$


$ $

\item \textbf{\underline{Multiplication}}

Compute the following polynomial multiplications in modulo $Q$: 

$d_0 = b^{\langle 1 \rangle}b^{\langle 2 \rangle} \bmod Q$

$d_1 = b^{\langle 2 \rangle}\vec{a}^{\langle 1 \rangle} +  b^{\langle 1 \rangle}\vec{a}^{\langle 2 \rangle} \bmod Q$

$\vec{d}_2 = \vec{a}^{\langle 1 \rangle} \otimes \vec{a}^{\langle 2 \rangle} \bmod Q$


$ $

\item \textbf{\underline{Relinearization}} 


Compute the following:

$ \textsf{ct}_\alpha = (d
_1, d_0)$

$ \textsf{ct}_\beta = \bm{\langle} \textsf{Decomp}^{\beta, l}(\vec{d}_2), \text{ } \textsf{Lev}_{\vec{s}, \sigma}^{\beta, l}(\vec{s}^{\otimes}) \bm{\rangle}$.

$ \textsf{ct}_{\alpha + \beta} = \textsf{ct}_{\alpha} + \textsf{ct}_{\beta}$

$ $

Then, the following property holds:

$\phi_{+}(\textsf{ct}_{\alpha+\beta})\equiv(\Delta m^{\langle1\rangle}+e^{\langle1\rangle}+qk_1)(\Delta m^{\langle2\rangle}+e^{\langle2\rangle}+qk_2)+e_{\rm rl}\pmod Q$, where $\phi_+(\vec a,b)=b+\langle\vec a,\vec s\rangle$.

$ $




\item \textbf{\underline{Rescaling}}

Update $\textsf{ct}_{\alpha + \beta} = (\vec{a}_{\alpha + \beta}, b_{\alpha + \beta}) \bmod Q$ to $\textsf{ct'}_{\alpha + \beta} = \left(\left\lceil\dfrac{\vec{a}_{\alpha + \beta}}{\Delta}\right\rfloor, \left\lceil\dfrac{b_{\alpha + \beta}}{\Delta}\right\rfloor\right) \bmod q$. 

$ $

This plaintext rescaling process can be also viewed as a modulus switch of the ciphertext $\textsf{ct}_{\alpha + \beta}$ from $Q \rightarrow q$. 


$ $

\end{enumerate}

Note that after the ciphertext-to-ciphertext multiplication, the plaintext scaling factor $\Delta=\dfrac{q}{t}$, the ciphertext modulus $q$, and the private key $\vec{s}$ stay the same as before.



\end{tcolorbox}



\textbf{Evaluation key and conditional correctness.} Let the LWE dimension be $k$, choose integer gadget weights $g_j=\beta^j$ for $0\leq j<L$, and require $\beta^L\geq Q$. For every ordered pair $(u,v)$, publish rows with phase $g_js_us_v+e_{u,v,j}^{\rm key}\pmod Q$. Decompose each component of $\vec d_2$ exactly modulo $Q$ into digits bounded by $D_g$. If $|e_{u,v,j}^{\rm key}|\leq B_{\rm key}$, the resulting relinearization error satisfies
\[
|e_{\rm rl}|\leq k^2LD_gB_{\rm key}.
\]
This key is distinct from TFHE's GGSW bootstrapping key. Use centered input components and messages, put $h=\|\vec s\|_1$, and suppose $|e_i|\leq B_i$. Then the integer phase carries obey
\[
|k_i|\leq\kappa_i:=\frac{h+2}{2}+\frac{B_i}{q}.
\]
If $R_a,R_b$ are the component rounding errors in the final division by $\Delta$, direct expansion gives
\[
\begin{aligned}
\phi_+(\textsf{ct}_{\rm out})&\equiv\Delta m_1m_2+\eta\pmod q,\\
\eta&=m_1e_2+m_2e_1+t(e_1k_2+e_2k_1)
      +\frac{e_1e_2+e_{\rm rl}}\Delta
      +\langle R_a,\vec s\rangle+R_b.
\end{aligned}
\]
The terms $qk_1m_2$, $qk_2m_1$ and $qt k_1k_2$ vanish modulo $q$. In particular, the dangerous modulus-remainder term of the general BFV teaching variant vanishes because $t\mid q$. A sufficient decoder condition is
\[
\frac t2(B_1+B_2)+t(B_1\kappa_2+B_2\kappa_1)
+\frac{B_1B_2+k^2LD_gB_{\rm key}}\Delta+\frac{h+1}{2}
<\frac\Delta2.
\]
Message reduction modulo $t$ adds a multiple of $q$ and no further error. This proves the displayed order conditionally. If one instead scales the integer tensor first by $1/\Delta=t/q$ and then relinearizes modulo $q$, the key modulus is $q$, the key-error term is no longer divided by $\Delta$, and tensor rounding is bounded by $(1+h+h^2)/2$. These are two separately justified orders, not interchangeable keys. The accompanying exact checks cover both orders with nonzero input and evaluation-key errors; they do not establish secure parameter sizes or a speed improvement.

\subsection{Homomorphic Key Switching}
\label{subsec:bfv-key-switching}

The gadget and its encrypted rows must use the same convention. For a general modulus $q$, one may use integer weights $\beta^{i-1}$ and enough digits for exact decomposition; a truncated decomposition adds its residual multiplied by the old secret to the key-switching error.

BFV's key switching scheme changes an RLWE ciphertext's secret key from $S$ to $S'$. This scheme is essentially RLWE's key switching scheme with the sign of the $A\cdot S$ term flipped in the encryption and decryption formula. Specifically, this is equivalent to the alternative GLWE version's (\autoref{subsec:glwe-alternative}) key switching scheme (\autoref{sec:glwe-key-switching}) with $k = 1$ as follows:

\begin{tcolorbox}[title={\textbf{\tboxlabel{\ref*{subsec:bfv-key-switching}} BFV's Key Switching}}]
$\textsf{RLWE}_{S',\sigma}(\Delta M) = (0, B) + \bm{\langle} \textsf{Decomp}^{\beta, l}(A), \text{ } \textsf{RLev}_{S', \sigma}^{\beta, l}(S) \bm{\rangle}$
\end{tcolorbox}




\subsection{Homomorphic Rotation of Input Vector Slots}
\label{subsec:bfv-rotation}


In this section, we will explain how to homomorphically rotate the elements of an input vector $\vec{v}$ after it is already encoded as a polynomial and encrypted as an RLWE ciphertext. In \autoref{subsec:coeff-rotation}, we learned how to rotate the coefficients of a polynomial. However, rotating the plaintext polynomial $M(X)$ or RLWE ciphertext polynomials $(A(X), B(X))$ does not necessarily rotate the input vector, which is the source of them. 

The key requirement of homomorphic rotation of input vector slots (i.e., input vector) is that this operation should be performed on the RLWE ciphertext such that after this operation, if we decrypt the RLWE ciphertext and decode it, the recovered input vector will be in a rotated state as we expect. We will divide this task into the following two sub-problems:

\begin{enumerate}[leftmargin=3\parindent]
\item How to indirectly rotate the input vector by updating the plaintext polynomial $M$ to $M'$?

\item How to indirectly update the plaintext polynomial $M$ to $M'$ by updating the RLWE ciphertext polynomials $(A, B)$ to $(A', B')$?
\end{enumerate}



\subsubsection{Rotating Input Vector Slots by Updating the Plaintext Polynomial}

In this task, our goal is to modify the plaintext polynomial $M(X)$ such that the first-half elements of the input vector $\vec{v}$ are shifted to the left by $h$ positions in a wrapping manner among them, and the second-half elements of $\vec{v}$ are also shifted to the left by $h$ positions in a wrapping manner among them. Specifically, if $\vec{v}$ is defined as follows:

$\vec{v} = (v_0, v_1, \cdots, v_{n-1})$

$ $

Then, we will denote the $h$-shifted vector $\vec{v}^{\langle h \rangle}$ as follows:  

$\vec{v}^{\langle h \rangle} = (\underbrace{v_h, v_{h+1}, \cdots, v_0, v_1, \cdots, v_{h-2}, v_{h-1},}_{\text{The first-half $\dfrac{n}{2}$ elements $h$-rotated to the left}} \text{ } \underbrace{v_{\frac{n}{2}+h}, v_{\frac{n}{2}+h+1}, \cdots, v_{\frac{n}{2}+h-2}, v_{\frac{n}{2}+h-1}}_{\text{The second-half $\dfrac{n}{2}$ elements $h$-rotated to the left}})$

$ $

Remember from \autoref{subsec:bfv-batch-encoding} that the BFV encoding scheme's components are as follows:

$\vec{v} = (v_0, v_1, v_2, \cdots, v_{n-1})$ \textcolor{red}{  $\rhd$ $n$-slot input vector}

$ $

$W =  \begin{bmatrix}
1 & 1 & 1 & \cdots & 1\\
(\omega) & (\omega^3) & (\omega^5) & \cdots & (\omega^{2n-1})\\
(\omega)^2 & (\omega^3)^2 & (\omega^5)^2 & \cdots & (\omega^{2n-1})^2\\
\vdots & \vdots & \vdots & \ddots & \vdots \\
(\omega)^{n-1} & (\omega^3)^{n-1} & (\omega^5)^{n-1} & \cdots & (\omega^{2n-1})^{n-1}\\
\end{bmatrix}$

$= \begin{bmatrix}
1 & 1 & \cdots & 1 & 1 & \cdots & 1 & 1\\
(\omega) & (\omega^3) & \cdots & (\omega^{n - 1}) & (\omega^{-(n - 1)}) & \cdots & (\omega^{-3}) & (\omega^{-1})\\
(\omega)^2 & (\omega^3)^2 & \cdots & (\omega^{n - 1})^2 & (\omega^{-(n - 1)})^2 & \cdots & (\omega^{-3})^2 & (\omega^{-1})^2\\
\vdots & \vdots & \vdots & \ddots & \vdots \\
(\omega)^{n-1} & (\omega^3)^{n-1} & \cdots & (\omega^{n - 1})^{n-1} & (\omega^{-(n - 1)})^{n-1} & \cdots & (\omega^{-3})^{n-1} & (\omega^{-1})^{n-1}\\
\end{bmatrix}$

, where $\omega = g^{\frac{t - 1}{2n}} \bmod t$ ($g$ is a generator of $\mathbb{Z}_t^{\times}$)

\textcolor{red}{ $\rhd$ The encoding matrix that converts $\vec{v}$ into $\vec{m}$ ( i.e., $n$ coefficients of the plaintext polynomial $M(X)$ )}
 


$ $

$\Delta\vec{m}  = n^{-1}\cdot\Delta W\cdot I_n^R\cdot\vec{v}$  

\textcolor{red}{  $\rhd$ A vector containing the  scaled $n$ integer coefficients of the plaintext polynomial}

$ $

$\Delta M = \sum\limits_{i=0}^{n-1} ( \Delta m_i  X^i)$ 

\textcolor{red}{ $\rhd$ The integer polynomial that isomorphically encodes the input vector $\vec{v}$}


$ $

We learned from \autoref{subsec:poly-vector-transformation} that decoding the polynomial $M(X)$ $\vec{v}$ is equivalent to evaluating $M(X)$ at the following $n$ distinct primitive $(\mu=2n)$-th root of unity: $\{\omega, \omega^3, \omega^5, \cdots, \omega^{2n-3}, \omega^{2n-1} \}$. Thus, the above decoding process is equivalent to the following:

$\vec{v} = \dfrac{W^T \Delta \vec{m}}{\Delta} = \bm{(}W^T_0\cdot \vec{m}, \text{ } W^T_1\cdot \vec{m}, \text{ } W^T_2\cdot \vec{m}, \text{ } \cdots, \text{ } W^T_{n-1}\cdot \vec{m}\bm{)}$  

$= \bm{(} \text{ } M(\omega), \text{ } M(\omega^3), \text{ } M(\omega^5), \cdots, M(\omega^{2n-3}), \text{ } M(\omega^{2n-1}) \text{ } \bm{)}$


{\small $= \bm{(} \text{ } M(\omega), \text{ } M(\omega^3), \text{ } M(\omega^5), \cdots, M(\omega^{n-3}), \text{ } M(\omega^{n-1}), \text{ } M(\omega^{-(n-1)}), \text{ } M(\omega^{-(n-3)}), \cdots, M(\omega^{-3}), \text{ } M(\omega^{-1}) \text{ } \bm{)}$}


%%$\textsf{LWE}_{S, \sigma}(\Delta M) = (A, B) \in \mathcal{R}_{\langle n, q \rangle }^{2}$ \textcolor{red}{  $\rhd$ An LWE ciphertext that encrypts $\Delta M$}

%$A\cdot S + B = \Delta M + E$


$ $

Now, our next task is to modify $M(X)$ to $M'(X)$ such that decoding $M'(X)$ will give us a modified input vector $\vec{v}^{\langle h \rangle}$ that is a rotation of the first half elements of $\vec{v}$ by $h$ positions to the left (in a wrapping manner among them), and the second half elements of it also rotated by $h$ positions to the left (in a wrapping manner among them). To accomplish this rotation, we will take a 2-step solution: 

\begin{enumerate}[leftmargin=3\parindent]
\item To convert $M(X)$ into $M'(X)$, we will define the new mapping $\sigma_M$ as follows: 

$\sigma_M : (M(X), h) \in (\mathcal{R}_{\langle n, t \rangle}, \mathbb{Z}_{n}) \longrightarrow M'(X) \in \mathcal{R}_{\langle n, t \rangle}$

, where $h$ is the number of rotation positions to be applied to $\vec{v}$.
\item To decode $M'(X)$ into the rotated input vector $\vec{v}^{\langle h \rangle}$, we need to re-design our decoding scheme by modifying \textsf{Encoding\textsubscript{1}}'s (\autoref{subsubsec:bfv-encoding-1}) isomorphic mapping $\sigma : M(X) \in \mathcal{R}_{\langle n, t \rangle} \longrightarrow \vec{v} \in \mathbb{Z}^n$ 
\end{enumerate}

$ $

\para{Converting $\bm{M(X)}$ into $\bm{M'(X)}$:} Our first task is to convert $M(X)$ into $M'(X)$, which is equivalent to applying our new mapping $\sigma_M : (M(X), h) \in (\mathcal{R}_{\langle n, t \rangle}, \mathbb{Z}_{n}) \longrightarrow M'(X) \in \mathcal{R}_{\langle n, t \rangle}$, such that decoding $M'(X)$ gives a rotated input vector $\vec{v}^{\langle h \rangle}$ whose first half of the elements in $\vec{v}$ are rotated by $h$ positions to the left (in a wrapping manner among them), and the second half of the elements are also rotated by $h$ positions to the left (in a wrapping manner among them). To design $\sigma_M$ that satisfies this requirement, we will use the number $5^j$ which has the following two special properties (based on number theory):
\begin{itemize}
\item $(5^j \bmod 2n)$ and $(-5^j \bmod 2n)$ generate all odd numbers between $[0, 2n)$ for the integer $j$ where $0 \leq j < \dfrac{n}{2}$. 
\item For each integer $j$ where $0 \leq j < \dfrac{n}{2}$, $(5^j \bmod 2n) + (-5^j \bmod 2n) \equiv 0 \bmod 2n$. 
\end{itemize}
For example, suppose the modulus $2n = 16$. Then, 

\begin{multicols}{2}
\noindent $5^0 \bmod 16 = 1$
\newline $5^1 \bmod 16 = 5$
\newline $5^2 \bmod 16 = 9$
\newline $5^3 \bmod 16 = 13$
\newline $-(5)^0 \bmod 16 = 15$
\newline $-(5)^1 \bmod 16 = 11$
\newline $-(5)^2 \bmod 16 = 7$
\newline $-(5)^3 \bmod 16 = 3$
\end{multicols}


As shown above, $0 \leq j < 4$ generate all odd numbers between $[0, 16)$. Also, for each $j$ in $0 \leq j < 4$, $(5^j \bmod 16) + (-5^j \bmod 16) = 16$. 

$ $

Let's define $J(h) = 5^h \bmod 2n$, \text{ } and \text{ } $J_*(h) = -5^h \bmod 2n$. Based on $J(h)$ and $J_*(h)$, we will define the mapping $\sigma_M : M(X) \rightarrow M'(X)$ as follows:

$\sigma_{M} : M(X) \rightarrow M(X^{J(h)})$

$ $

Given a plaintext polynomial $M(X)$, in order to give its decoded version of input vector $\vec{v}$ the effect of the first half of the elements being rotated by $h$ positions to the left (in a wrapping manner) and the second half of the elements also being rotated by $h$ positions to the left (in a wrapping manner), we update the current plaintext polynomial $M(X)$ to a new polynomial $M'(X) = M(X^{J(h)}) = M(X^{5^h})$ by applying the $\sigma_M$ mapping, where $h$ is the number of positions for left rotations for the first half and second half of the elements of $\vec{v}$. 


$ $

\para{Decoding $\bm{M'(X)}$ into \boldmath$\vec{v}^{\langle h \rangle}$}: Our second task is to modify our original decoding scheme in order to successfully decode $M'(X)$ into the rotated input vector $\vec{v}^{\langle h \rangle}$. For this, we will modify our original isomorphic mapping $\sigma : M(X) \longrightarrow \vec{v}$, from:

$\sigma: M(X) \in \mathcal{R}_{\langle n, q \rangle} \longrightarrow  \bm{(}M(\omega), M(\omega^3),M(\omega^5), \cdots, M(\omega^{2n-1})\bm{)} \in \mathbb{Z}^n$ \textcolor{red}{ \text{ } \# designed in \autoref{subsec:poly-vector-transformation}}

$ $

, to the following: 

$\sigma_J: M(X) \in \mathcal{R}_{\langle n, t \rangle} \longrightarrow  \bm{(}M(\omega^{J(0)}),M(\omega^{J(1)}),M(\omega^{J(2)}), \cdots, M(\omega^{J(\frac{n}{2} - 1)}), $

\textcolor{white}{$\sigma_J: M(X) \in \mathcal{R}_{\langle n, q \rangle} \longrightarrow  \bm{(}$} $M(\omega^{J_*(0)}), \text{ } M(\omega^{J_*(1)}), \text{ } M(\omega^{J_*(2)}), \cdots, M(\omega^{J_*(\frac{n}{2} - 1)}) \text{ } \bm{)} \in \mathbb{Z}^{n}$


$ $

The common aspect between $\sigma$ and $\sigma_J$ is that they both evaluate the polynomial $M(X)$ at $n$ distinct primitive $(\mu=2n)$-th roots of unity (i.e., $w^i$ for all odd $i$ between $[0,2n]$ ). In the case of the $\sigma_J$ mapping, note that $J(j) = 5^j \bmod 2n$ \text{ } and \text{ } $J_*(j) = -5^j \bmod 2n$ \text{ } for each $j$ in $0 \leq j < \dfrac{n}{2} $ \text{ }  cover all odd numbers between $[0, 2n]$. Therefore, $\omega^{J(j)}$ and $\omega^{J_*(j)}$ between $0 \leq j < \frac{n}{2}$ cover all $n$ distinct primitive $(\mu=2n)$-th roots of unity.

Meanwhile, the difference between $\sigma$ and $\sigma_J$ is the order of the output vector elements. In the $\sigma$ mapping, the order of evaluated coordinates for $M(X)$ is $\omega, \omega^3, \cdots, \omega^{2n - 1}$, whereas in the $\sigma_J$ mapping, the order of evaluated coordinates is $\omega^{J(0)}, \omega^{J(1)}, \cdots, \omega^{J(\frac{n}{2} - 1)}, \omega^{J_*(0)}, \omega^{J_*(1)}, \cdots, \omega^{J_*(\frac{n}{2} - 1)}$. We will later explain why we modified the ordering like this. 

In the original \textsf{Decoding\textsubscript{2}} process (\autoref{subsec:bfv-batch-encoding}), applying the $\sigma$ mapping to a plaintext polynomial $M(X)$ was equivalent to computing the following: 

$\vec{v}_{'} = \bm{(} \text{ } M(\omega), \text{ } M(\omega^3), \text{ } M(\omega^5), \cdots, M(\omega^{2n-3}), \text{ } M(\omega^{2n-1}) \bm{)}$ 

$ $

$ = \bm{(} \text{ } M(\omega), \text{ } M(\omega^3), \text{ } M(\omega^5), \cdots, M(\omega^{n-1}), \text{ } M(\omega^{-(n-1)}), \cdots, \text{ } M(\omega^{-3}), \text{ } M(\omega^{-1}) \text{ } \bm{)}$ 

$ $

$= \bm{(} \text{ } W^T_0\cdot \vec{m}, \text{ } W^T_1\cdot \vec{m}, \text{ } W^T_2\cdot \vec{m}, \text{ } \cdots, \text{ } W^T_{n-1}\cdot \vec{m} \text{ } \bm{)} $ 
$ $ \textcolor{red}{ $\rhd$ where $W_i^T$ is the $(i+1)$-th row of $W^T$}

$ = W^T \vec{m} $
%$= \dfrac{I_n^R}{n}\cdot \bm{(} \text{ } M(\omega), \text{ } M(\omega^3), \text{ } M(\omega^5), \cdots, M(\omega^{2n-3}), \text{ } M(\omega^{2n-1}) \text{ } )$




$ $

Similarly, the modified $\sigma_J$ mapping to the plaintext polynomial $M(X)$ is equivalent to computing the following:

$\vec{v}_{'} =  \bm{(} \text{ } 
M(\omega^{J(0)}), \text{ } M(\omega^{J(1)}), \text{ } M(\omega^{J(2)}), \cdots,  M(\omega^{J(\frac{n}{2}-1)}), \text{ } M(\omega^{J_*(0)}), \text{ } M(\omega^{J_*(1)}), \cdots,  M(\omega^{J_*(\frac{n}{2}-1)}) \text{ } \bm{)}$


$ $

$= \bm{(} \text{ } \hathat{W}^*_0\cdot \vec{m}, \text{ } \hathat{W}^*_1\cdot \vec{m}, \text{ } \hathat{W}^*_2\cdot \vec{m}, \text{ } \cdots, \text{ } \hathat{W}^*_{n-1}\cdot \vec{m} \text{ } \bm{)} $ 

$ $

$ = \hathat{W}^* \vec{m} $, where 

$ $

\noindent $\hathat{W}^* = \begin{bmatrix}
1 & (\omega^{J(0)}) & (\omega^{J(0)})^2 & \cdots & (\omega^{J(0)})^{n-1}\\
1 & (\omega^{J(1)}) & (\omega^{J(1)})^2 & \cdots & (\omega^{J(1)})^{n-1}\\
1 & (\omega^{J(2)}) & (\omega^{J(2)})^2 & \cdots & (\omega^{J(2)})^{n-1}\\
\vdots & \vdots & \vdots & \ddots & \vdots \\
1 & (\omega^{J(\frac{n}{2}-1)}) & (\omega^{J(\frac{n}{2}-1)})^2 & \cdots & (\omega^{J(\frac{n}{2}-1)})^{n-1}\\
1 & (\omega^{J_*(0)}) & (\omega^{J_*(0)})^2 & \cdots & (\omega^{J_*(0)})^{n-1}\\
1 & (\omega^{J_*(1)}) & (\omega^{J_*(1)})^2 & \cdots & (\omega^{J_*(1)})^{n-1}\\
1 & (\omega^{J_*(2)}) & (\omega^{J_*(2)})^2 & \cdots & (\omega^{J_*(2)})^{n-1}\\
\vdots & \vdots & \vdots & \ddots & \vdots \\
1 & (\omega^{J_*(\frac{n}{2}-1)}) & (\omega^{J_*(\frac{n}{2}-1)})^2 & \cdots & (\omega^{J_*(\frac{n}{2}-1)})^{n-1}\\
\end{bmatrix}$

$ $

$ $

\noindent $\hathat{W} = \begin{bmatrix}
1 & 1 & \cdots & 1 & 1 & 1 & \cdots & 1\\
(\omega^{J(\frac{n}{2} - 1)}) & (\omega^{J(\frac{n}{2} - 2)}) & \cdots & (\omega^{J(0)}) & (\omega^{J_*(\frac{n}{2} - 1)}) & (\omega^{J_*(\frac{n}{2} - 2)}) & \cdots & (\omega^{J_*(0)})\\
(\omega^{J(\frac{n}{2} - 1)})^2 & (\omega^{J(\frac{n}{2} - 2)})^2 & \cdots & (\omega^{J(0)})^2 & (\omega^{J_*(\frac{n}{2} - 1)})^2 & (\omega^{J_*(\frac{n}{2} - 2)})^2 & \cdots & (\omega^{J_*(0)})^2 \\
\vdots & \vdots & \ddots & \vdots & \vdots & \vdots & \vdots & \vdots \\
(\omega^{J(\frac{n}{2} - 1)})^{n-1} & (\omega^{J(\frac{n}{2} - 2)})^{n-1} & \cdots & (\omega^{J(0)})^{n-1} & (\omega^{J_*(\frac{n}{2} - 1)})^{n-1} & (\omega^{J_*(\frac{n}{2} - 2)})^{n-1} & \cdots  & (\omega^{J_*(0)})^{n-1}
\end{bmatrix}$

$ $

$ $

From this point, we will replace $W$ in the \textsf{Encoding\textsubscript{1}} process (\autoref{subsec:bfv-batch-encoding}) by $\hathat{W}$, and $W^T$ in the \textsf{Decoding\textsubscript{2}} process by $\hathat{W}^*$. 

$ $

To demonstrate that $\hathat{W}$ is a valid encoding matrix like $W$ and $\hathat{W}^*$ is a valid decoding matrix like $W^T$, we need to prove the following 2 aspects:
\begin{itemize}[leftmargin=3\parindent]

\item \textbf{${\hathat{\bm W}}$ is a basis of the $\bm{n}$-dimensional vector space:} This is true, because $\hathat{W}$ is simply a column-wise re-ordering of $W$, which is still a basis of the $\bm{n}$-dimensional vector space.

\item \boldmath{$\hathat{W}^{*} \cdot \hathat{W} = n \cdot I_n^R$ (to satisfy Theorem~\ref*{subsec:vandermonde-euler} in \autoref{subsec:vandermonde-euler})}: This proof is split into 2 sub-proofs: 
\begin{enumerate}[leftmargin=2\parindent]
\item \boldmath{${\hathat{ W}^{ *} \cdot \hathat{ W}}$ has value $\bm{n}$ along the anti-diagonal line:} Each element along the anti-diagonal line of \boldmath$\hathat{ W}^{ *} \cdot \hathat{ W}$ is computed as $\sum\limits_{i=0}^{n-1} \omega^{2nik} = n$ where $k$ is some integer.
\item \boldmath{${\hathat{ W}^{ *} \cdot \hathat{ W}}$ has value 0 at all other coordinates:} All other elements except for the ones along the anti-diagonal lines are $\sum\limits_{i=0}^{n-1}r^i=\dfrac{r^n-1}{r-1}=0\quad(r=\omega^{a+b}\ne1)$ (by Geometric Sum).
\end{enumerate}

We provide \href{https://github.com/wang2000star/fhe-textbook-cn/blob/chinese-edition/source%20code/bfv_j_matrix_inverse_proof.py}{\underline{the Python script}} that empirically demonstrates this.

\end{itemize}




Therefore, $\hathat{W}^*$ and $\hathat{W}$ are valid encoding \& decoding matrices that transform $\vec{v}$ into $M(X)$. 

$ $

Now, let's think about what will be the structure of $\vec{v}^{\langle h \rangle}$ (i.e., the first-half elements of $\vec{v}$ being rotated $h$ positions to the left in a wrapping manner among them and the second-half elements of it also being rotated $h$ positions to the left in a wrapping manner among them). Remember that $\vec{v}$ is as follows:

$\vec{v} = \bm{(} \text{ } 
M(\omega^{J(0)}), \text{ } M(\omega^{J(1)}), \cdots,  M(\omega^{J(\frac{n}{2}-1)}), \text{ } 
M(\omega^{J_*(0)}), \text{ } M(\omega^{J_*(1)}),  \cdots,  M(\omega^{J_*(\frac{n}{2}-1)}) \bm{)}$


\text{ } \text{ } $= \bm{(} \text{ } \hathat{W}^*_0\cdot \vec{m}, \text{ } \hathat{W}^*_1\cdot \vec{m}, \text{ } \hathat{W}^*_2\cdot \vec{m}, \text{ } \cdots, \text{ } \hathat{W}^*_{n - 1}\cdot \vec{m} \text{ } \bm{)} $ 

$ $

Thus, the state of $\vec{v}^{\langle h \rangle}$ which is equivalent to rotating $\vec{v}$ by $h$ positions to the left for the first-half and second-half element groups will be the following: 

$\vec{v}^{\langle h \rangle} = \bm{(} \text{ } M(\omega^{J(h)}), \text{ } M(\omega^{J(h+1)}), \cdots, M(\omega^{J(\frac{n}{2}-1)}), \text{ } M(\omega^{J(0)}), \text{ }M(\omega^{J(1)}), \cdots, \text{ } M(\omega^{J(h-2)}), \text{ } M(\omega^{J(h-1)}),$

\text{ }\text{ }\text{ } $ M(\omega^{J_*(h)}), \text{ } M(\omega^{J_*(h+1)}), \cdots, M(\omega^{J_*(\frac{n}{2}-1)}), \text{ } M(\omega^{J_*(0)}), \text{ }M(\omega^{J_*(1)}), \cdots, \text{ } M(\omega^{J_*(h-2)}), \text{ } M(\omega^{J_*(h-1)}) \text{ } \bm{)}$

$ $

Notice that the above computation of $\vec{v}^{\langle h \rangle}$ is equivalent to vertically rotating the upper $\frac{n}{2}$ rows of $\hathat{W}^*$ by $h$ positions upward (in a wrapping manner among them), rotating the lower $\frac{n}{2}$ rows of $\hathat{W}^{*}$ by $h$ positions upward (in a wrapping manner among them), and multiplying the resulting matrix with $\vec{m}$. However, it is not desirable to directly modify the decoding matrix $\hathat{W}^*$ like this in practice, because then the decoding matrix loses its consistency. Therefore, instead of directly modifying $\hathat{W}^*$, we will modify $\vec{m}$ to $\vec{m}_{'}$ (i.e., modify $M(X)$ to $M'(X)$) such that the relation $\vec{v}^{\langle h \rangle} = \hathat{W}^* \cdot \vec{m}_{'}$ holds. Let's extract the upper-half rows of $\hathat{W}^*$ and denote this $\dfrac{n}{2} \times n$ matrix as $\hathat{H}_1^*$. Then, $\hathat{H}_1^*$ is equivalent to a Vandermonde matrix (Definition~\ref*{subsec:vandermonde} in \autoref{subsec:vandermonde}) in the form of $V(\omega^{J(0)}, \omega^{J(1)}, \cdots, \omega^{J(\frac{n}{2}-1)})$. Similarly, let's extract the lower-half rows of $\hathat{W}^*$ and denote this $\dfrac{n}{2} \times n$ matrix as $\hathat{H}_2^*$. Then, $\hathat{H}_2^*$ is equivalent to a Vandermonde matrix (Definition~\ref*{subsec:vandermonde} in \autoref{subsec:vandermonde}) in the form of $V(\omega^{J_*(0)}, \omega^{J_*(1)}, \cdots, \omega^{J_*(\frac{n}{2}-1)})$. 

Now, let's vertically rotate the rows of $\hathat{H}_1^*$ by $h$ positions upward and denote it as $\hathat{H}_1^{*\langle h \rangle}$; and vertically rotate the rows of $\hathat{H}_2^*$ by $h$ positions upward and denote it as $\hathat{H}_2^{*\langle h \rangle}$. And let's denote the $\dfrac{n}{2}$-dimensional vector comprising the first-half elements of $\vec{v}^{\langle h \rangle}$ as $\vec{v}_1^{\langle h \rangle}$, and the $\dfrac{n}{2}$-dimensional vector comprising the second-half elements of $\vec{v}^{\langle h \rangle}$ as $\vec{v}_2^{\langle h \rangle}$. 
Then, computing (i.e., decoding) $\vec{v}_1^{\langle h \rangle} = \hathat{H}_1^{*\langle h \rangle}\cdot \vec{m}$ is equivalent to modifying $M(X)$ to $M'(X) = M(X^{J(h)})$ (whose coefficient vector is $\vec{m}^{\langle h \rangle}$) and then computing (i.e., decoding) $\vec{v}_1^{\langle h \rangle} = \hathat{H}_1^{*}\cdot \vec{m}^{\langle h \rangle}$. This is because:

$\vec{v}_1^{\langle h \rangle} =\hathat{H}_1^{*\langle h \rangle}\cdot \vec{m}$

$ = \bm{(} \text{ } \hathat{W}^*_{h}\cdot\vec{m}, \text{ } \hathat{W}^*_{h+1}\cdot\vec{m}, \text{ } \hathat{W}^*_{h+2}\cdot\vec{m}, \cdots, \hathat{W}^*_{\frac{n}{2}-1}\cdot\vec{m}, \text{ } \hathat{W}^*_{0}\cdot\vec{m}, \text{ }\hathat{W}^*_{1}\cdot\vec{m}, \cdots, \text{ } \hathat{W}^*_{h-2}\cdot\vec{m}, \text{ } \hathat{W}^*_{h-1}\cdot\vec{m}  \bm{)}$

$ $

$= \bm{(} \text{ } M((\omega^{J(h)})^{J(0)}), \text{ } M((\omega^{J(h)})^{J(1)}), \text{ } M((\omega^{J(h)})^{J(2)}), \cdots, M((\omega^{J(h)})^{J(\frac{n}{2}-1)}) \text{ } \bm{)}$


\textcolor{red}{ $\rhd$ This is equivalent to evaluating the polynomial $M(X^{J(h)})$ at the following $\dfrac{n}{2}$ distinct $(\mu=2n)$-th roots of unity: $\omega^{J(0)}, \omega^{J(1)},  \omega^{J(2)}, \cdots, \omega^{J(\frac{n}{2} - 1)}$}

$ $

$= \bm{(} \text{ } M(\omega^{J(h)\cdot J(0)}), \text{ } M(\omega^{J(h)\cdot J(1)}), \text{ } M(\omega^{J(h)\cdot J(2)}), \cdots, M(\omega^{J(h)\cdot J(\frac{n}{2}-1)}) \text{ } \bm{)}$

$= \bm{(} \text{ } M(\omega^{5^h\cdot 5^0}), \text{ } M(\omega^{5^h\cdot 5^1}), \text{ } M(\omega^{5^h\cdot 5^2}), \cdots, M(\omega^{5^h\cdot 5^{n/2-1}}) \text{ } \bm{)}$

$= \bm{(} \text{ } M(\omega^{5^{h}}), \text{ } M(\omega^{5^{h+1}}), \text{ } M(\omega^{5^{h+2}}), \cdots, M(\omega^{5^{h+n/2-1}}) \text{ } \bm{)}$ \textcolor{red}{ $\rhd$ note that $5^{\frac{n}{2}} \bmod 2n = 1$}

$ $

$= \bm{(} \text{ } M(\omega^{J(h)}), \text{ } M(\omega^{J(h+1)}), \text{ } M(\omega^{J(h+2)}), \cdots, M(\omega^{J(\frac{n}{2}-1)}), \text{ } M(\omega^{J(0)}), \text{ }M(\omega^{J(1)}),$ 
                                    
$\cdots, \text{ } M(\omega^{J(h-2)}), \text{ } M(\omega^{J(h-1)})  \bm{)}$ 

$ $

$= \hathat{H}_1^{*}\cdot \vec{m}^{\langle h \rangle}$ \textcolor{red}{ \text{ } \# where $\vec{m}^{\langle h \rangle}$ contains the $n$ coefficients of the polynomial $M(X^{J(h)})$ }


$ $

Similarly, computing (i.e., decoding) $\vec{v}_2^{\langle h \rangle} = \hathat{H}_2^{*\langle h \rangle}\cdot \vec{m}$ is equivalent to modifying $M(X)$ to $M'(X) = M(X^{J(h)})$ (whose coefficient vector is $\vec{m}^{\langle h \rangle}$) and then computing (i.e., decoding) $\vec{v}_2^{\langle h \rangle} = \hathat{H}_2^{*}\cdot \vec{m}^{\langle h \rangle}$. This is because,

$\vec{v}_2^{\langle h \rangle} =\hathat{H}_2^{*\langle h \rangle}\cdot \vec{m}$

$ $


$ = \bm{(} \text{ } \hathat{W}^*_{\frac{n}{2} +h}\cdot\vec{m}, \text{ } \hathat{W}^*_{\frac{n}{2} + h+1}\cdot\vec{m}, \text{ } \hathat{W}^*_{\frac{n}{2} +h+2}\cdot\vec{m}, \cdots, \hathat{W}^*_{n -1}\cdot\vec{m}, \text{ } \hathat{W}^*_{\frac{n}{2}}\cdot\vec{m}, \text{ }\hathat{W}^*_{\frac{n}{2} + 1}\cdot\vec{m}, \cdots,$

\text{ }\text{ }\text{ }\text{ }$ \text{ } \hathat{W}^*_{\frac{n}{2} + h-2}\cdot\vec{m}, \text{ } \hathat{W}^*_{\frac{n}{2} + h-1}\cdot\vec{m}  \bm{)}$

$ $

$= \bm{(} \text{ } M((\omega^{J(h)})^{J_*(0)}), \text{ } M((\omega^{J(h)})^{J_*(1)}), \text{ } M((\omega^{J(h)})^{J_*(2)}), \cdots, M((\omega^{J(h)})^{J_*(\frac{n}{2}-1)}) \text{ } \bm{)}$


\textcolor{red}{ $\rhd$ This is equivalent to evaluating the polynomial $M(X^{J(h)})$ at the following $\dfrac{n}{2}$ distinct $(\mu=2n)$-th roots of unity: $\omega^{J_*(0)}, \omega^{J_*(1)},  \omega^{J_*(2)}, \cdots, \omega^{J_*(\frac{n}{2} - 1)}$}

$ $

$= \bm{(} \text{ } M(\omega^{J(h)\cdot J_*(0)}), \text{ } M(\omega^{J(h)\cdot J_*(1)}), \text{ } M(\omega^{J(h)\cdot J_*(2)}), \cdots, M(\omega^{J(h)\cdot J_*(\frac{n}{2}-1)}) \text{ } \bm{)}$

$= \bm{(} \text{ } M(\omega^{5^h\cdot -5^0}), \text{ } M(\omega^{5^h\cdot -5^1}), \text{ } M(\omega^{5^h\cdot -5^2}), \cdots, M(\omega^{5^h\cdot -5^{n/2-1}}) \text{ } \bm{)}$

$= \bm{(} \text{ } M(\omega^{-5^{h}}), \text{ } M(\omega^{-5^{h+1}}), \text{ } M(\omega^{-5^{h+2}}), \cdots, M(\omega^{-5^{h+n/2-1}}) \text{ } \bm{)}$ \textcolor{red}{ $\rhd$ note that $-(5^{\frac{n}{2}} \bmod 2n) = -1$}

$ $

$= \bm{(} \text{ } M(\omega^{J_*(h)}), \text{ } M(\omega^{J_*(h+1)}), \text{ } M(\omega^{J_*(h+2)}), \cdots, M(\omega^{J_*(\frac{n}{2}-1)}), \text{ } M(\omega^{J_*(0)}), \text{ }M(\omega^{J_*(1)}),$ 
                                    
$\cdots, \text{ } M(\omega^{J_*(h-2)}), \text{ } M(\omega^{J_*(h-1)})  \bm{)}$ 

$ $

$= \hathat{H}_2^{*}\cdot \vec{m}^{\langle h \rangle}$


$ $

The above derivations demonstrate that $\vec{v}_1^{\langle h \rangle} = \hathat{H}_1^{*}\cdot \vec{m}^{\langle h \rangle}$, and $\vec{v}_2^{\langle h \rangle} = \hathat{H}_2^{*}\cdot \vec{m}^{\langle h \rangle}$. Combining these two findings, we reach the conclusion that $\vec{v}^{\langle h \rangle} = \hathat{W}^* \cdot \vec{m}^{\langle h \rangle}$: rotating the first-half elements of the input vector $\vec{v}$ by $h$ positions to the left and the second-half elements of it by $h$ positions also to the left is equivalent to updating the plaintext polynomial $M(X)$ to $M(X^{J(h)})$ and then decoding it with the decoding matrix $\hathat{W}^*$. 



However, now a new problem is that we cannot directly update the plaintext $M(X)$ to $M(X^{J(h)})$, because $M(X)$ is encrypted as an RLWE ciphertext. Therefore, we need to instead update the RLWE ciphertext components $(A, B)$ to \textit{indirectly} by updating $M(X)$ to $M(X^{J(h)})$. We will explain this in the next subsection.


\subsubsection{Updating the Plaintext Polynomial by Updating the Ciphertext Polynomials}

Given an RLWE ciphertext $\textsf{ct} = (A, B)$, our goal is to update $\textsf{ct} = (A, B)$ to $C^{\langle h \rangle} = (A^{\langle h \rangle}, B^{\langle h \rangle})$ such that decrypting it gives the input vector $\vec{v}^{\langle h \rangle}$. That is, the following relation should hold: 

$\textsf{RLWE}^{-1}_{S, \sigma}\bm{(} \text{ } C^{\langle h \rangle} = (A^{\langle h \rangle}, B^{\langle h \rangle}) \text{ }\bm{)} = \Delta M(X^{J(h)}) + E'$

$ $

Remember that in the RLWE cryptosystem (\autoref{sec:rlwe})'s alternative version (\autoref{subsec:glwe-alternative}), the plaintext and ciphertext pair have the following relation:

$\Delta M(X) + E(X) = A(X)\cdot S(X) + B(X) \approx \Delta M(X)$

$ $

If we apply $X = X^{J(h)}$ in the above relation, we can derive the following relation: 

$\Delta M(X^{J(h)}) + E(X^{J(h)}) = A(X^{J(h)}) \cdot S(X^{J(h)}) + B(X^{J(h)}) \approx \Delta M(X^{J(h)})$

$ $

This relation implies that if we decrypt the ciphertext $C^{\langle h \rangle} = (A(X^{J(h)}), B(X^{J(h)}))$ with $S(X^{J(h)})$ as the secret key, then we get $\Delta M(X^{J(h)})$. Therefore, $C^{\langle h \rangle} = (A(X^{J(h)}), B(X^{J(h)}))$ is the RLWE ciphertext we are looking for, because decrypting it and then decoding its plaintext $M(X^{J(h)})$ will give us the input vector $\vec{v}^{\langle h \rangle}$. 

Apply the automorphism componentwise to obtain $C^{\langle h\rangle}=(A(X^{J(h)}),B(X^{J(h)}))$, whose secret is $S(X^{J(h)})$. Then key-switch to the original secret $S(X)$ using Summary~\ref*{subsec:bfv-key-switching} from \autoref{subsec:bfv-key-switching}. This produces new ciphertext components and adds key-switching error; the same unchanged pair is not simultaneously an encryption under both secrets. The following expression suppresses that error only in the message notation:

$ $

\noindent $\underbrace{\textsf{RLWE}_{S(X),\sigma}\bm{(}\Delta M(X^{J(h)})\bm{)}}_{\substack{\text{the result of} \\ \text{plaintext-to-ciphertext addition}}} = \underbrace{\bm{(} \text{ } 0, B(X^{J(h)}) \text{ } \bm{)}}_{\substack{\text{the plaintext } B(X^{J(h)}) \\ \text{(trivial ciphertext)}}} + \bm{\langle} \text{ } \underbrace{\textsf{Decomp}^{\beta, l}\bm{(}A(X^{J(h)})\bm{)}, \text{ } \textsf{RLev}_{S(X), \sigma}^{\beta, l}\bm{(}S(X^{J(h)})\bm{)} \text{ } \bm{\rangle}}_{\substack{\substack{\text{an RLWE ciphertext encrypting $A(X^{J(h)}) \cdot S(X^{J(h)})$}\\ \text{which is key-switched from $S(X^{J(h)})\rightarrow S(X)$}}}}$



\subsubsection{Summary}
\label{subsubsec:bfv-rotation-summary}

We summarize the procedure of rotating the BFV input vectors as follows: 

\begin{tcolorbox}[title={\textbf{\tboxlabel{\ref*{subsec:bfv-rotation}} BFV's Homomorphic Rotation of Input Vector Slots}}]

To support input vector slot rotation, we update the original encoding matrix in \textsf{Encoding\textsubscript{1}} as follows: 

$ $

{\footnotesize{\noindent $\hathat{W} = \begin{bmatrix}
1 & 1 & \cdots & 1 & 1 & 1 & \cdots & 1\\
(\omega^{J(\frac{n}{2} - 1)}) & (\omega^{J(\frac{n}{2} - 2)}) & \cdots & (\omega^{J(0)}) & (\omega^{J_*(\frac{n}{2} - 1)}) & (\omega^{J_*(\frac{n}{2} - 2)}) & \cdots & (\omega^{J_*(0)})\\
(\omega^{J(\frac{n}{2} - 1)})^2 & (\omega^{J(\frac{n}{2} - 2)})^2 & \cdots & (\omega^{J(0)})^2 & (\omega^{J_*(\frac{n}{2} - 1)})^2 & (\omega^{J_*(\frac{n}{2} - 2)})^2 & \cdots & (\omega^{J_*(0)})^2 \\
\vdots & \vdots & \ddots & \vdots & \vdots & \vdots & \ddots & \vdots \\
(\omega^{J(\frac{n}{2} - 1)})^{n-1} & (\omega^{J(\frac{n}{2} - 2)})^{n-1} & \cdots & (\omega^{J(0)})^{n-1} & (\omega^{J_*(\frac{n}{2} - 1)})^{n-1} & (\omega^{J_*(\frac{n}{2} - 2)})^{n-1} & \cdots  & (\omega^{J_*(0)})^{n-1}
\end{bmatrix}$}}

$ $

$ $


, and update the original decoding matrix in \textsf{Decoding\textsubscript{2}} as follows:

$\hathat{W}^* = \begin{bmatrix}
1 & (\omega^{J(0)}) & (\omega^{J(0)})^2 & \cdots & (\omega^{J(0)})^{n-1}\\
1 & (\omega^{J(1)}) & (\omega^{J(1)})^2 & \cdots & (\omega^{J(1)})^{n-1}\\
1 & (\omega^{J(2)}) & (\omega^{J(2)})^2 & \cdots & (\omega^{J(2)})^{n-1}\\
\vdots & \vdots & \vdots & \ddots & \vdots \\
1 & (\omega^{J(\frac{n}{2}-1)}) & (\omega^{J(\frac{n}{2}-1)})^2 & \cdots & (\omega^{J(\frac{n}{2}-1)})^{n-1}\\
1 & (\omega^{J_*(0)}) & (\omega^{J_*(0)})^2 & \cdots & (\omega^{J_*(0)})^{n-1}\\
1 & (\omega^{J_*(1)}) & (\omega^{J_*(1)})^2 & \cdots & (\omega^{J_*(1)})^{n-1}\\
1 & (\omega^{J_*(2)}) & (\omega^{J_*(2)})^2 & \cdots & (\omega^{J_*(2)})^{n-1}\\
\vdots & \vdots & \vdots & \ddots & \vdots \\
1 & (\omega^{J_*(\frac{n}{2}-1)}) & (\omega^{J_*(\frac{n}{2}-1)})^2 & \cdots & (\omega^{J_*(\frac{n}{2}-1)})^{n-1}\\
\end{bmatrix}$

$ $

$ $

, where $J(h)$ is the \textit{rotation helper formula}: $J(h) = 5^h \bmod 2n$, 
\text{ } $J_*(h) = -5^h \bmod 2n$. 

$ $

Using $\hathat{W}$, the encoding is perform as: $\vec{m} = n^{-1}\cdot \hathat{W} \cdot I_n^R \cdot \vec{v}$. Using $\hathat{W}^*$, the decoding is performed as $\vec{v} = \hathat{W}^* \cdot \vec{m}$

$ $

Suppose we have an RLWE ciphertext and a key-switching key as follows:

$\textsf{RLWE}_{S, \sigma}(\Delta M) = (A, B)$, \text{ } $\textsf{RLev}_{S, \sigma}^{\beta, l}\bm(S(X^{J(h)})\bm)$

$ $

Then, the procedure of rotating the first-half elements of the ciphertext's original input vector $\vec{v}$ by $h$ positions to the left (in a wrapping manner among them) and the second-half elements of $\vec{v}$ by $h$ positions to the left (in a wrapping manner among them) is as follows: 

\begin{enumerate}
\item Update $A(X)$, $B(X)$ to $A(X^{J(h)})$, $B(X^{J(h)})$. 
\item Perform the following key switching (\autoref{subsec:bfv-key-switching}) from $S(X^{J(h)})$ to $S(X)$:

$\textsf{RLWE}_{S(X),\sigma}\bm{(}\Delta M(X^{J(h)})\bm{)} = \bm{(} 0, B(X^{J(h)}) \bm{)} \text{ } + \text{ } \bm{\langle}  \textsf{Decomp}^{\beta, l}\bm{(}A(X^{J(h)})\bm{)}, \text{ } \textsf{RLev}_{S(X), \sigma}^{\beta, l}\bm{(}S(X^{J(h)})\bm{)} \bm{\rangle}$
\end{enumerate}


\end{tcolorbox}


\subsubsection{Encoding Example}
\label{subsubsec:bfv-batch-encoding-ex}

%$bfv_vector = [10, 0, 5, 13]
%bfv_vector2 = [2, 4, 8, 6]


In the above example, we use the unsigned modulo representation (e.g., $[0, t-1], [0, q]$) when computing $\bmod t$ and $\bmod q$. However, the correctness of the result holds even if we use the centered modulo representation (e.g., $[-\dfrac{t}{2}, \dfrac{t}{2})$, $[-\dfrac{q}{2}, \dfrac{q}{2})$). In many actual FHE libraries, centered modulo arithmetic is often used to manage noise growth more effectively.

Suppose we have the following setup: 

$\mu = 8, \, n = \dfrac{\mu}{2} = 4, \, t = 17, \, q = 2^8 = 256, \, n^{-1} = 13, \, \Delta = \left\lfloor\dfrac{q}{t}\right\rfloor = 15$

$ $

$\mathcal{R}_{\langle 4, 17 \rangle} = \mathbb{Z}_{17}[X] / (X^4 + 1)$

$ $

The roots of $X^4 + 1 \pmod{17}$ are $X = \{2, 8, 15, 9\}$, as demonstrated as follows:

$2^4 \equiv 8^4 \equiv 15^4 \equiv 9^4 \equiv 16 \equiv -1 \bmod{17}$

$ $

Definition~\ref*{subsec:cyclotomic-def} (in \autoref{subsec:cyclotomic-def}) states that the roots of the $\mu$-th cyclotomic polynomial are the primitive $\mu$-th roots of unity. And Definition~\ref*{subsec:roots-def} (in \autoref{subsec:roots-def}) states that the order of the primitive $\mu$-th roots of unity is $\mu$. These definitions apply to both the cyclotomic polynomials over $X\in\mathbb{C}$ (complex numbers) and the cyclotomic polynomials over $X \in \mathbb{Z}_t$ (ring). 

Since $\{2, 8, 15, 9\}$ are the roots of the $(\mu=8)$-th cyclotomic polynomial $X^4 + 1$ over the ring $\mathbb{Z}_{17}$, they are also the $(\mu=8)$-th primitive roots of unity of $\mathbb{Z}_{17}$. Therefore, their order (\autoref{subsec:order-def}) is $\mu=8$ as demonstrated as follows: 

$2^8 \equiv 8^8 \equiv 15^8 \equiv 9^8 \equiv 1 \bmod 17$

$2^4 \equiv 8^4 \equiv 15^4 \equiv 9^4 \equiv 16 \not\equiv 1 \bmod 17$

$ $

Definition~\ref*{subsec:cyclotomic-def} (in \autoref{subsec:cyclotomic-def}) and Theorem~\ref*{subsec:cyclotomic-theorem} (\autoref{subsec:cyclotomic-theorem}) also state that for each primitive $\mu$-th root of unity $\omega$, $\{\omega^k\}_{\textsf{gcd}(k, \mu) = 1}$ generates all roots of the $\mu$-th cyclotomic polynomial. Notice that in the case of the $(\mu=8)$-th cyclotomic polynomial $X^4 + 1$, its roots $\{2, 8, 15, 9\}$ generate all $(\mu=8)$-th roots of unity as follows:

$\{2^1, 2^3, 2^5, 2^7\} \equiv \{8^1, 8^3, 8^5, 8^7\} \equiv \{15^1, 15^3, 15^5, 15^7\} \equiv \{9^1, 9^3, 9^5, 9^7\} \equiv \{2, 8, 15, 9\} \bmod 17$

$ $

Among $\{2, 8, 15, 9\}$ as the roots of $X^4 + 1$, let's choose $\omega = 9$ as the base root to construct the encoding matrix $\hathat{W}$ and the decoding matrix $\hathat{W}^*$ as follows: 

$\hathat{W} = \begin{bmatrix}
1 & 1 & 1 & 1\\
\omega^{J(1)} & \omega^{J(0)} & \omega^{J_*(1)} & \omega^{J_*(0)}\\
(\omega^{J(1)})^2 & (\omega^{J(0)})^2 & (\omega^{J_*(1)})^2 & (\omega^{J_*(0)})^2\\
(\omega^{J(1)})^3 & (\omega^{J(0)})^3 & (\omega^{J_*(1)})^3 & (\omega^{J_*(0)})^3\\
\end{bmatrix}$ \textcolor{red}{  $\rhd$ where $J(h) = 5^h \bmod 8$}

$ = \begin{bmatrix}
1 & 1 & 1 & 1\\
9^{5} & 9^{1} & 9^{3} & 9^{7}\\
(9^{5})^2 & (9^{1})^2 & (9^{3})^2 & (9^{7})^2\\
(9^{5})^3 & (9^{1})^3 & (9^{3})^3 & (9^{7})^3\\
\end{bmatrix} \equiv \begin{bmatrix}
1 & 1 & 1 & 1\\
8 & 9 & 15 & 2\\
13 & 13 & 4 & 4\\
2 & 15 & 9 & 8\\
\end{bmatrix} \bmod{17}$ 

$ $ 

$\hathat{W}^* = \begin{bmatrix}
1 & \omega^{J(0)} & (\omega^{J(0)})^2 & (\omega^{J(0)})^3\\
1 & \omega^{J(1)} & (\omega^{J(1)})^2 & (\omega^{J(1)})^3\\
1 & \omega^{J_*(0)} & (\omega^{J_*(0)})^2 & (\omega^{J_*(0)})^3\\
1 & \omega^{J_*(1)} & (\omega^{J_*(1)})^2 & (\omega^{J_*(1)})^3\\
\end{bmatrix} \equiv \begin{bmatrix}
1 & 9 & 13 & 15\\
1 & 8 & 13 & 2\\
1 & 2 & 4 & 8\\
1 & 15 & 4 & 9\\
\end{bmatrix}  \bmod{17}$

$ $

Notice that Theorem~\ref*{subsec:vandermonde-euler-integer-ring} (in \autoref{subsec:vandermonde-euler-integer-ring}) is demonstrated as follows:

$\hathat{W}^* \cdot \hathat{W} = \begin{bmatrix}
1 & 9 & 13 & 15\\
1 & 8 & 13 & 2\\
1 & 2 & 4 & 8\\
1 & 15 & 4 & 9\\
\end{bmatrix} \cdot \begin{bmatrix}
1 & 1 & 1 & 1\\
8 & 9 & 15 & 2\\
13 & 13 & 4 & 4\\
2 & 15 & 9 & 8\\
\end{bmatrix} = \begin{bmatrix}
0 & 0 & 0 & 4\\
0 & 0 & 4 & 0\\
0 & 4 & 0 & 0\\
4 & 0 & 0 & 0\\
\end{bmatrix} = n I_n^{R} \pmod{17}$

$ $

Now suppose that we encode the following two input vectors (i.e., input vector slots) in $\mathbb{Z}_{17}$:

$\vec{v}_1 = (10, 3, 5, 13)$

$\vec{v}_2 = (2, 4, 3, 6)$

$\vec{v}_1 + \vec{v}_2 = (10, 3, 5, 13) + (2, 4, 3, 6) \equiv (12, 7, 8, 2) \bmod 17$

$ $

These two vectors are encoded as follows:

$ $

$\vec{m}_1 = n^{-1} \hathat{W} \cdot I_n^R \cdot \vec{v} = 13 \cdot \begin{bmatrix}
1 & 1 & 1 & 1\\
8 & 9 & 15 & 2\\
13 & 13 & 4 & 4\\
2 & 15 & 9 & 8\\
\end{bmatrix} \cdot \begin{bmatrix}
0 & 0 & 0 & 1\\
0 & 0 & 1 & 0\\
0 & 1 & 0 & 0\\
1 & 0 & 0 & 0\\
\end{bmatrix} \cdot \begin{bmatrix}
10\\
3\\
5\\
13\\
\end{bmatrix} \equiv (12, 11, 12, 1)  \bmod 17$

$ $

$\vec{m}_2 = n^{-1} \hathat{W} \cdot I_n^R \cdot \vec{v} = 13 \cdot \begin{bmatrix}
1 & 1 & 1 & 1\\
8 & 9 & 15 & 2\\
13 & 13 & 4 & 4\\
2 & 15 & 9 & 8\\
\end{bmatrix} \cdot \begin{bmatrix}
0 & 0 & 0 & 1\\
0 & 0 & 1 & 0\\
0 & 1 & 0 & 0\\
1 & 0 & 0 & 0\\
\end{bmatrix} \cdot \begin{bmatrix}
2\\
4\\
3\\
6\\
\end{bmatrix} \equiv (8, 5, 14, 6) \bmod 17$

$ $

$ $

$\Delta M_1(X) = 15\cdot(12 + 11X + 12X^2 + 1X^3) = 180 + 165X + 180X^2 + 15X^3 \pmod{q}$ \textcolor{red}{  $\rhd$ where $q = 256$}

$\Delta M_2(X) = 15\cdot(8 + 5X + 14X^2 + 6X^3) = 120 + 75X + 210X^2 + 90X^3 \pmod{q}$

$\Delta M_{1+2}(X) = \Delta M_1(X) + \Delta M_2(X) = \Delta (M_1(X) + M_2(X)) = 44 + 240X + 134X^2 + 105X^3 \pmod{q}$

$M_{1+2}(X) =\left\lceil\dfrac{17}{256}(44 + 240X + 134X^2 + 105X^3)\right\rfloor \bmod 17$

$= \left\lceil2.922 + 15.938X + 8.898X^2 + 6.973X^3\right\rfloor \bmod 17$

$3 + 16X + 9X^2 + 7X^3 \bmod 17 $

$ $


This polynomial matches the value of $\vec{m}_{1+2}$ as follows:

$\vec{m}_1 + \vec{m}_2 = (12, 11, 12, 1) + (8, 5, 14, 6) = (20, 16, 26, 7) \equiv (3, 16, 9, 7) \bmod 17$

$ $

Finally, we decode $\vec{m}_{1+2}$ as follows:



$\vec{v}_{1+2} = \hathat{W}^* \cdot \vec{m}_{1+2} = \begin{bmatrix}
1 & 9 & 13 & 15\\
1 & 8 & 13 & 2\\
1 & 2 & 4 & 8\\
1 & 15 & 4 & 9\\
\end{bmatrix} \cdot \begin{bmatrix}3\\16\\9\\7\end{bmatrix} = (12, 7, 8, 2) \bmod{17}$

$ $

This result matches the expected vector sum $\vec{v}_1 + \vec{v}_2$.




\subsubsection{Rotation Example}
\label{subsubsec:bfv-rotation-ex}

Also in the above example, we use the unsigned modulo representation (e.g., $[0, t-1], [0, q]$) when computing $\bmod t$ and $\bmod q$. However, the correctness of the result holds even if we use the centered modulo representation (e.g., $[-\dfrac{t}{2}, \dfrac{t}{2})$, $[-\dfrac{q}{2}, \dfrac{q}{2})$).

Suppose we have the following setup: 

$\mu = 16, n = \dfrac{\mu}{2} = 8, \text{ } t = 17, \text{ } q = 2^8 = 256, \text{ } n^{-1} = 15, \text{ } \Delta = 15$

$ $

$\mathcal{R}_{\langle 8, 17 \rangle} = \mathbb{Z}_{17}[X] / (X^8 + 1)$

$ $

The roots of $X^8 + 1 \pmod{17}$ are $X = \{3, 5, 6, 7, 10, 11, 12, 14\}$, as demonstrated as follows:

$3^8 \equiv 5^8 \equiv 6^8 \equiv 7^8 \equiv 10^8 \equiv 11^8 \equiv 12^8 \equiv 14^8 \equiv 16 \equiv -1 \bmod{17}$


$ $

Among $\{3, 5, 6, 7, 10, 11, 12, 14\}$ as the roots of $X^8 + 1$, let's choose $\omega = 3$ as the base root to construct the encoding matrix $\hathat{W}$ and the decoding matrix $\hathat{W}^*$ as follows: 

$\hathat{W} = \begin{bmatrix}
1 & 1 & 1 & 1 & 1 & 1 & 1 & 1\\
\omega^{J(3)} & \omega^{J(2)} & \omega^{J(1)} & \omega^{J(0)} & \omega^{J_*(3)} & \omega^{J_*(2)} & \omega^{J_*(1)} & \omega^{J_*(0)}\\
(\omega^{J(3)})^2 & (\omega^{J(2)})^2 & (\omega^{J(1)})^2 & (\omega^{J(0)})^2 & (\omega^{J_*(3)})^2 & (\omega^{J_*(2)})^2 & (\omega^{J_*(1)})^2 & (\omega^{J_*(0)})^2\\
(\omega^{J(3)})^3 & (\omega^{J(2)})^3 & (\omega^{J(1)})^3 & (\omega^{J(0)})^3 & (\omega^{J_*(3)})^3 & (\omega^{J_*(2)})^3 & (\omega^{J_*(1)})^3 & (\omega^{J_*(0)})^3\\
(\omega^{J(3)})^4 & (\omega^{J(2)})^4 & (\omega^{J(1)})^4 & (\omega^{J(0)})^4 & (\omega^{J_*(3)})^4 & (\omega^{J_*(2)})^4 & (\omega^{J_*(1)})^4 & (\omega^{J_*(0)})^4\\
(\omega^{J(3)})^5 & (\omega^{J(2)})^5 & (\omega^{J(1)})^5 & (\omega^{J(0)})^5 & (\omega^{J_*(3)})^5 & (\omega^{J_*(2)})^5 & (\omega^{J_*(1)})^5 & (\omega^{J_*(0)})^5\\
(\omega^{J(3)})^6 & (\omega^{J(2)})^6 & (\omega^{J(1)})^6 & (\omega^{J(0)})^6 & (\omega^{J_*(3)})^6 & (\omega^{J_*(2)})^6 & (\omega^{J_*(1)})^6 & (\omega^{J_*(0)})^6\\
(\omega^{J(3)})^7 & (\omega^{J(2)})^7 & (\omega^{J(1)})^7 & (\omega^{J(0)})^7 & (\omega^{J_*(3)})^7 & (\omega^{J_*(2)})^7 & (\omega^{J_*(1)})^7 & (\omega^{J_*(0)})^7\\
\end{bmatrix}$ 

$ \equiv \begin{bmatrix}
1 & 1 & 1 & 1 & 1 & 1 & 1 & 1\\
12&14&5&3&10&11&7&6\\
8&9&8&9&15&2&15&2\\
11&7&6&10&14&5&3&12\\
13&13&13&13&4&4&4&4\\
3&12&14&5&6&10&11&7\\
2&15&2&15&9&8&9&8\\
7&6&10&11&5&3&12&14\\
\end{bmatrix} \bmod{17}$ 

$ $ 

$\hathat{W}^* = \begin{bmatrix}
1 & \omega^{J(0)} & (\omega^{J(0)})^2 & (\omega^{J(0)})^3 & (\omega^{J(0)})^4 & (\omega^{J(0)})^5 & (\omega^{J(0)})^6 & (\omega^{J(0)})^7\\
1 & \omega^{J(1)} & (\omega^{J(1)})^2 & (\omega^{J(1)})^3 & (\omega^{J(1)})^4& (\omega^{J(1)})^5& (\omega^{J(1)})^6& (\omega^{J(1)})^7\\
1 & \omega^{J(2)} & (\omega^{J(2)})^2 & (\omega^{J(2)})^3 & (\omega^{J(2)})^4 & (\omega^{J(2)})^5 & (\omega^{J(2)})^6 & (\omega^{J(2)})^7\\
1 & \omega^{J(3)} & (\omega^{J(3)})^2 & (\omega^{J(3)})^3 & (\omega^{J(3)})^4& (\omega^{J(3)})^5& (\omega^{J(3)})^6& (\omega^{J(3)})^7\\
1 & \omega^{J_*(0)} & (\omega^{J_*(0)})^2 & (\omega^{J_*(0)})^3 & (\omega^{J_*(0)})^4 & (\omega^{J_*(0)})^5 & (\omega^{J_*(0)})^6 & (\omega^{J_*(0)})^7\\
1 & \omega^{J_*(1)} & (\omega^{J_*(1)})^2 & (\omega^{J_*(1)})^3 & (\omega^{J_*(1)})^4& (\omega^{J_*(1)})^5& (\omega^{J_*(1)})^6& (\omega^{J_*(1)})^7\\
1 & \omega^{J_*(2)} & (\omega^{J_*(2)})^2 & (\omega^{J_*(2)})^3 & (\omega^{J_*(2)})^4 & (\omega^{J_*(2)})^5 & (\omega^{J_*(2)})^6 & (\omega^{J_*(2)})^7\\
1 & \omega^{J_*(3)} & (\omega^{J_*(3)})^2 & (\omega^{J_*(3)})^3 & (\omega^{J_*(3)})^4& (\omega^{J_*(3)})^5& (\omega^{J_*(3)})^6& (\omega^{J_*(3)})^7\\
\end{bmatrix}$

$ \equiv \begin{bmatrix}
1&3&9&10&13&5&15&11\\
1&5&8&6&13&14&2&10\\
1&14&9&7&13&12&15&6\\
1&12&8&11&13&3&2&7\\
1&6&2&12&4&7&8&14\\
1&7&15&3&4&11&9&12\\
1&11&2&5&4&10&8&3\\
1&10&15&14&4&6&9&5\\
\end{bmatrix}  \bmod{17}$

$ $

$ $

Notice that Theorem~\ref*{subsec:vandermonde-euler-integer-ring} (in \autoref{subsec:vandermonde-euler-integer-ring}) is demonstrated as follows:

$\hathat{W}^* \cdot \hathat{W} = \begin{bmatrix}
1&3&9&10&13&5&15&11\\
1&5&8&6&13&14&2&10\\
1&14&9&7&13&12&15&6\\
1&12&8&11&13&3&2&7\\
1&6&2&12&4&7&8&14\\
1&7&15&3&4&11&9&12\\
1&11&2&5&4&10&8&3\\
1&10&15&14&4&6&9&5\\
\end{bmatrix} \cdot \begin{bmatrix}
1 & 1 & 1 & 1 & 1 & 1 & 1 & 1\\
12&14&5&3&10&11&7&6\\
8&9&8&9&15&2&15&2\\
11&7&6&10&14&5&3&12\\
13&13&13&13&4&4&4&4\\
3&12&14&5&6&10&11&7\\
2&15&2&15&9&8&9&8\\
7&6&10&11&5&3&12&14\\
\end{bmatrix}$

$ = \begin{bmatrix}
0 & 0 & 0 & 0 & 0 & 0 & 0 & 8\\
0 & 0 & 0 & 0 & 0 & 0 & 8 & 0\\
0 & 0 & 0 & 0 & 0 & 8 & 0 & 0\\
0 & 0 & 0 & 0 & 8 & 0 & 0 & 0\\
0 & 0 & 0 & 8 & 0 & 0 & 0 & 0\\
0 & 0 & 8 & 0 & 0 & 0 & 0 & 0\\
0 & 8 & 0 & 0 & 0 & 0 & 0 & 0\\
8 & 0 & 0 & 0 & 0 & 0 & 0 & 0\\
\end{bmatrix} = n I_n^{R} \pmod{17}$

$ $

Now suppose that we encode the following input vector (i.e., input vector slots) in $\mathbb{Z}_{17}$:

$\vec{v} = (1, 2, 3, 4, 5, 6, 7, 8)$

By rotating this vector by 3 positions to the left (i.e., the first-half slots and the second-half slots separately wrapping around within their own group), we get a new vector:

$\vec{v}_r = (4, 1, 2, 3, 8, 5, 6, 7) \bmod 17$

$ $

$\vec{v}$ is encoded as follows:

$ $

$\vec{m} = n^{-1} \hathat{W} \cdot I_n^R \cdot \vec{v} $

$= 15 \cdot  \begin{bmatrix}
1 & 1 & 1 & 1 & 1 & 1 & 1 & 1\\
12&14&5&3&10&11&7&6\\
8&9&8&9&15&2&15&2\\
11&7&6&10&14&5&3&12\\
13&13&13&13&4&4&4&4\\
3&12&14&5&6&10&11&7\\
2&15&2&15&9&8&9&8\\
7&6&10&11&5&3&12&14\\
\end{bmatrix} \cdot \begin{bmatrix}
0 & 0 & 0 & 0 & 0 & 0 & 0 & 1\\
0 & 0 & 0 & 0 & 0 & 0 & 1 & 0\\
0 & 0 & 0 & 0 & 0 & 1 & 0 & 0\\
0 & 0 & 0 & 0 & 1 & 0 & 0 & 0\\
0 & 0 & 0 & 1 & 0 & 0 & 0 & 0\\
0 & 0 & 1 & 0 & 0 & 0 & 0 & 0\\
0 & 1 & 0 & 0 & 0 & 0 & 0 & 0\\
1 & 0 & 0 & 0 & 0 & 0 & 0 & 0\\
\end{bmatrix} \cdot \begin{bmatrix}
1\\
2\\
3\\
4\\
5\\
6\\
7\\
8\\
\end{bmatrix}$

$ \equiv (13, 16, 10, 5, 9, 12, 7, 1)  \bmod 17$


$ $

$\Delta M(X) = 15\cdot(13 + 16X + 10X^2 + 5X^3 + 9X^4 + 12X^5 + 7X^6 + X^7)$


$ = 195 + 240X + 150X^2 + 75X^3 + 135X^4 + 180X^5 + 105X^6 + 15X^7 \pmod{q}$ \textcolor{red}{  $\rhd$ where $q = 256$}

$ $

This polynomial matches the value of $\Delta M(X^{J(3)})$ as follows:

$\Delta M(X^{J(3)}) = \Delta M(X^{13}) $

$= 195 + 180X - 150X^2 - 15X^3 + 135X^4 - 240X^5 - 105X^6 + 75X^7 \pmod{q}$

$ $



%Note that the coefficients of the scaled polynomials $\Delta M(X^{J(3)})$ are still within the range of the ciphertext domain $q=64$ (which must hold throughout all homomorphic computations to preserve correctness). 

Now, we decode $M(X^{J(3)})$ as follows: 

$\vec{m}_{J(3)} = \dfrac{\Delta \vec{m}_{J(3)}}{\Delta} = \dfrac{(195, 180, -150, -15, 135, -240, -105, 75)}{15} = (13, 12, -10, -1, 9, -16, -7, 5) \pmod{17}$

$ $


$\vec{v}_{J(3)} = \hathat{W}^* \cdot \vec{m}_{J(3)} = \begin{bmatrix}
1&3&9&10&13&5&15&11\\
1&5&8&6&13&14&2&10\\
1&14&9&7&13&12&15&6\\
1&12&8&11&13&3&2&7\\
1&6&2&12&4&7&8&14\\
1&7&15&3&4&11&9&12\\
1&11&2&5&4&10&8&3\\
1&10&15&14&4&6&9&5\\
\end{bmatrix} \cdot \begin{bmatrix}
13\\
12\\
-10\\
-1\\
9\\
-16\\
-7\\
5\end{bmatrix}$

$= (4, 1, 2, 3, 8, 5, 6, 7) \pmod{17}$

$= \vec{v}_r$

$ $

The decoded $\vec{v}_{J(3)}$ matches the expected rotated input vector $\vec{v}_r$.  


$ $

In practice, apply the automorphism to the encrypted components, producing
\[
\textsf{ct}^{\langle h=3\rangle}=(A(X^{J(3)}),B(X^{J(3)}))
\]
. Its decryption secret is $S(X^{J(3)})$. Switch back to $S(X)$ with the corresponding rotation key before ordinary decryption; provided the transformed and key-switching errors remain within the bound, decoding yields $\vec v_r$.

$ $




\para{Source Code:} Examples of BFV's batch encoding and homomorphic input vector rotation can be executed by running \href{https://github.com/wang2000star/fhe-textbook-cn/blob/chinese-edition/source%20code/bfv.py}{\underline{this Python script}} (BFV addition) and \href{https://github.com/wang2000star/fhe-textbook-cn/blob/chinese-edition/source%20code/bfv_rotation_only.py}{\underline{this one}} (BFV rotation only).



\subsection{Application: Matrix Multiplication}
\label{subsec:bfv-matrix-multiplication}

BFV supports packed inner products by multiplying slots and summing them with rotations and additions. A doubling schedule uses logarithmically many rotations within a power-of-two slot block; summing both rows additionally needs a row exchange. For matrix-vector multiplication, the diagonal method below assumes $m\leq n$ and a logical cyclic rotation $\rho$ of all $n$ active slots. The previous subsection implements two independent half-row rotations, so it cannot by itself realize this full $\rho$. Either place the logical vector in one cyclic row of a larger ring, or combine row rotations, the row-swap automorphism $X\mapsto X^{-1}$ and plaintext masks. Given a matrix $A$ and vector $\vec x$ (Definition~\ref*{subsec:matrix-arithmetic} in \autoref{subsec:matrix-arithmetic}):

$A = \begin{bmatrix}
 a_{\langle 0, 0\rangle} & a_{\langle 0, 1\rangle} & a_{\langle 0, 2\rangle} & \cdots & a_{\langle 0, n-1\rangle}\\
 a_{\langle 1, 0\rangle} & a_{\langle 1, 1\rangle} & a_{\langle 1, 2\rangle} & \cdots & a_{\langle 1, n-1\rangle} \\
 a_{\langle 2, 0\rangle} & a_{\langle 2, 1\rangle} & a_{\langle 2, 2\rangle} & \cdots & a_{\langle 2, n-1\rangle} \\
\vdots & \vdots & \vdots & \ddots & \vdots \\
 a_{\langle m-1, 0\rangle} & a_{\langle m-1, 1\rangle} & a_{\langle m-1, 2\rangle} & \cdots & a_{\langle m-1, n-1\rangle} \\
\end{bmatrix} = \begin{bmatrix} 
\vec{a}_{\langle 0, * \rangle} \\ 
\vec{a}_{\langle 1, * \rangle} \\ 
\vec{a}_{\langle 2, * \rangle} \\ 
\vdots\\
\vec{a}_{\langle m-1, * \rangle} \\ 
\end{bmatrix}, \text{ } \vec{x} = (x_0, x_1, \cdots, x_{n-1})$

$ $

The result of $A \cdot \vec{x}$ is an $m$-dimensional vector computed as:

$A \cdot \vec{x} = \Big(\vec{a}_{\langle 0, * \rangle} \cdot \vec{x}, \text{ } \vec{a}_{\langle 1, * \rangle} \cdot \vec{x}, \text{ }\cdots,\text{ } \vec{a}_{\langle m-1, * \rangle} \cdot \vec{x} \Big) = \left(\sum\limits_{i=0}^{n-1} a_{0, i} \cdot  x_i, \text{ } \sum\limits_{i=0}^{n-1}  a_{1, i} \cdot x_i, \cdots, \text{ } \sum\limits_{i=0}^{n-1} a_{m-1, i} \cdot x_i \right)$

$ $

Let's define $\rho(\vec{v}, h)$ as the rotation of $\vec{v}$ by $h$ positions to the left. And remember that the Hadamard dot product (Definition~\ref*{subsec:vector-arithmetic} in \autoref{subsec:vector-arithmetic}) is defined as slot-wise multiplication of two vectors: 

$\vec{a} \odot \vec{b} = (a_0 b_0, \text{ } a_1 b_1, \text{ } \cdots, \text{ } a_{n-1} b_{n-1})$

$ $

Let's define $n$ distinct diagonal vector $\vec{u}_i$ extracted from matrix $A$ as follows: 

$\vec{u}_i = \{a_{\langle j \bmod m, \text{ } (i + j)\bmod n \rangle}\}_{j = 0}^{n-1}$

$ $

Then, the original matrix-to-vector multiplication formula can be equivalently constructed as follows:

$A\cdot\vec x=\left(\sum_{i=0}^{n-1}\vec u_i\odot\rho(\vec x,i)\right)_{0:m}$

, whose computation result is equivalent to $A \cdot \vec{x}$. 
The above formula is compatible with homomorphic computation, because BFV supports Hadamard dot product between input vectors as a ciphertext-to-plaintext multiplication between their polynomial-encoded forms (\autoref{subsec:bfv-mult-plain}), and BFV also supports $\rho(\vec{v}, h)$ as homomorphic input vector slot rotation (\autoref{subsec:bfv-rotation}). After homomorphically computing the above formula, we can consider only the first $m$ (out of $n$) resulting input vector slots to store the result of $A \cdot \vec{x}$. 



\subsection{Noise Bootstrapping}
\label{subsec:bfv-bootstrapping}

\noindent \textbf{- Reference 1:} 
\href{https://eprint.iacr.org/2022/1363.pdf}{Bootstrapping for BGV and BFV Revisited}~\cite{cryptoeprint:2022/1363}


\noindent \textbf{- Reference 2:} 
\href{https://eprint.iacr.org/2014/873.pdf}{Bootstrapping for HELib}~\cite{10.1007/s00145-020-09368-7}

\noindent \textbf{- Reference 3:} 
\href{https://arxiv.org/pdf/1906.02867}{A Note on Lower Digits Extraction Polynomial for Bootstrapping}~\cite{huo2019notelowerdigitsextraction}

\noindent \textbf{- Reference 4:} 
\href{https://eprint.iacr.org/2024/1587.pdf}{Fully Homomorphic Encryption for Cyclotomic Prime Moduli}~\cite{cryptoeprint:2024/1587}

%\noindent \textbf{- Reference 4:} 
%\href{https://s-space.snu.ac.kr/bitstream/10371/152893/1/000000155273.pdf}{Bootstrapping Methods for Homomorphic Encryption}

%\noindent \textbf{- Reference 5:} 
%\href{https://www.esat.kuleuven.be/cosic/publications/thesis-418.pdf}{Bootstrapping Algorithms for BGV and FV}

$ $

In BFV, continuous ciphertext-to-ciphertext multiplication increases the noise in a multiplicative manner, and once the noise overflows the message bits, then the message gets corrupted. Bootstrapping is a process of resetting the grown noise. 

\subsubsection{High-level Idea}
\label{subsubsec:bfv-bootstrapping-high-level}

Let $t=p^r$, $\varepsilon>r$ and $P=p^\varepsilon$. We decompose the decryption function, give a general coefficient/slot construction, and state an explicit gate-by-gate sufficient correctness test in \autoref{subsubsec:bfv-bootstrapping-summary}. The homomorphic polynomial evaluations in that test use standard FV integer tensor multiplication followed by $T/q$ scaling at the current plaintext modulus $T$, not the $1/\Delta$ teaching variant when $T\nmid q$. This is a conditional correctness result, not a secure or optimized parameter selection. Encrypting a secret under itself additionally needs the evaluation-key and circular-security assumptions of the cited construction~\cite{cryptoeprint:2022/1363}.

For a current plaintext modulus $T$, it is useful to measure error against the exact rational scale $\mu_T=q/T$. With $r_T=q-T\lfloor q/T\rfloor$, the same ciphertext satisfies
\[
\phi_S(C)\equiv\mu_TM_T+\mathcal E_T\pmod q,
\qquad \mathcal E_T=E_T-\frac{r_T}{T}M_T.
\]
Here $M_T$ is the centered coefficient representative. The decoder is correct if $\|\mathcal E_T\|_\infty<q/(2T)$. This change of error coordinates neither changes the ciphertext nor introduces an approximate plaintext scheme; it is exactly the decoder condition from \autoref{subsec:bfv-enc-dec}.

For a ciphertext $(A,B)$ with phase $AS+B=\Delta M+E\pmod q$, the target plaintext identity is
\[
A'=\Bigl[\Bigl\lceil(P/q)A\Bigr\rfloor\Bigr]_P,\quad
B'=\Bigl[\Bigl\lceil(P/q)B\Bigr\rfloor\Bigr]_P,\quad
Z=[A'S+B']_P,
\qquad M=\Bigl[\Bigl\lceil Z/p^{\varepsilon-r}\Bigr\rfloor\Bigr]_{p^r}.
\]
\begin{enumerate}
\item \textbf{Modulus Switch:} Change the input ciphertext modulus from $q$ to $P$. This alone does not change the original plaintext modulus $p^r$. For suitable parameters the switched phase is $p^{\varepsilon-r}M+E'\pmod P$, with $\|E'\|_\infty<p^{\varepsilon-r}/2$.
\item \textbf{Homomorphic Decryption:} Use an evaluation ciphertext encrypting $S$ at a large ciphertext modulus $q$, now with plaintext modulus $P$ and scale $\Delta'=\lfloor q/P\rfloor$. Compute $A'\cdot\textsf{RLWE}_{S,\sigma}(\Delta'S)+(0,\Delta'B')$. The result encrypts the phase $Z\pmod P$, with additional evaluation noise. The public constant $B'$ must be encoded at scale $\Delta'$.
\item \textbf{\textsf{CoeffToSlot}:} Apply the homomorphic linear transformation that places the coefficients of $Z$ into slots. Under the full scalar-splitting conditions of \autoref{subsubsec:bfv-rotation-summary}, the transformation on the slot vector is $n^{-1}\hathat W I_n^R$.
\item \textbf{Digit Extraction:} Evaluate polynomial low-digit extraction and exact-division steps to compute $[\lceil z_i/p^{\varepsilon-r}\rfloor]_{p^r}$ in each slot. For odd $p$, use balanced digits. For $p=2$, unsigned digit deletion is flooring, so a half-unit offset is required before deletion. The digit identities are explained below; their use in a complete bootstrap is subject to the CRT-domain restriction of \autoref{subsubsec:bfv-bootstrapping-coefftoslot}.
\item \textbf{\textsf{SlotToCoeff}:} Apply the inverse linear transformation $\hathat W^*$ to put the recovered coefficients back into a polynomial. The result encrypts $M$ with a new error $E^{\langle b\rangle}$ generated by the evaluation. Its bound depends on the parameters and circuit; it is not automatically smaller than an arbitrary input error $E$ or the differently scaled $E'$.
\end{enumerate}

The old $E'$ is removed as part of evaluating the rounding function, while the errors of the homomorphic evaluation remain. A usable bootstrap must leave enough output noise budget for subsequent operations. Repetition supports arbitrary prescribed finite circuit depth only with suitable parameters and an acceptable total failure probability.

\subsubsection{Modulus Switch}
\label{subsubsec:bfv-bootstrapping-modulus-switch}

Let $P=p^\varepsilon$ and define the component rounding polynomials $R_A=A'-(P/q)A$ and $R_B=B'-(P/q)B$, using compatible integer lifts before reduction. Their coefficient magnitudes are at most $1/2$. The exact switched phase is
\[
A'S+B'\equiv p^{\varepsilon-r}M+E'\pmod P,
\quad E'=\frac PqE+\left(\frac Pq\Delta-p^{\varepsilon-r}\right)M+R_AS+R_B.
\]
Thus a sufficient condition for the target rounding is $\|E'\|_\infty<p^{\varepsilon-r}/2$. The component-rounding term cannot be replaced by rounding the old error alone. Its deterministic bound is $(\|S\|_1+1)/2$, in addition to the scaled old error and the scale mismatch term.

\subsubsection{Homomorphic Decryption}
\label{subsubsec:bfv-bootstrapping-homomorphic-decryption}

Write $Z=[A'S+B']_P=[p^{\varepsilon-r}M+E']_P$. With $\Delta'=\lfloor q/P\rfloor$, let the evaluation ciphertext $(A_s,B_s)$ satisfy
\[
A_sS+B_s=\Delta'S+E_s\pmod q.
\]
The homomorphic inner product is the ciphertext
\[
C=(A'A_s,\ A'B_s+\Delta'B'),
\qquad \phi_S(C)=\Delta'(A'S+B')+A'E_s\pmod q.
\]
This explicitly distinguishes the encrypted message $\Delta'S$ from the encryption secret $S$. If $A'S+B'=Z+PK$, then
\[
\phi_S(C)\equiv\Delta'Z+A'E_s-r_PK\pmod q,
\qquad r_P=q-P\Delta'.
\]
The plaintext reduction modulo $P$ is therefore represented by an additional error $-r_PK$ when $P\nmid q$; it does not require a separate ciphertext-side modulo gate. Also $\|A'E_s\|_\infty\leq\|A'\|_1\|E_s\|_\infty$. The loose inequality $q\gg P$ alone does not ensure that this error fits a decoding budget of order $q/P$.

The desired coefficient operation is $[\lceil z_i/p^{\varepsilon-r}\rfloor]_{p^r}$. It is implemented by a sequence of polynomial evaluations and exact divisions, not by one degree-$p$ polynomial for the entire rounding function. To evaluate it coefficientwise, first perform \textsf{CoeffToSlot}, then the digit-extraction procedure of \autoref{subsubsec:bfv-bootstrapping-digit-extraction}, and finally \textsf{SlotToCoeff}. The homomorphic evaluation errors remain separate from the old error $E'$ inside the encrypted integer $z_i$.

\subsubsection{\textsf{CoeffToSlot} and \textsf{SlotToCoeff}}
\label{subsubsec:bfv-bootstrapping-coefftoslot}

The displayed $n$-scalar-slot matrices assume $p\equiv1\pmod{2n}$, with a primitive root Hensel-lifted to $\mathbb Z_{p^\varepsilon}$ and compatibly reduced at each later plaintext modulus. In general, when $\gcd(p,2n)=1$ and $d=\operatorname{ord}_{2n}(p)$, there are $n/d$ extension-ring slots. The trace-dual construction below packs all coefficients into scalar entries of $d$ ciphertexts, so complete scalar splitting is not necessary. Matrix multiplication denotes a homomorphic linear circuit using automorphisms, key switching and plaintext multiplications. This CRT route still requires $p$ odd in this chapter's power-of-two ring: at $p=2$, $X^n+1=(X+1)^n\pmod2$ has repeated factors. A binary CRT bootstrap must instead use a cyclotomic index coprime to $2$, or a separately justified construction.

\textbf{A general realizable linear transform.} To include the alternative binary ring, temporarily write $R_j=(\mathbb Z/p^j\mathbb Z)[X]/\Phi_{m_c}(X)$ with $p\nmid m_c$, $n=\varphi(m_c)$, $d=\operatorname{ord}_{m_c}(p)$ and $\ell=n/d$. This chapter's ring has $m_c=2n$. Choose compatible monic factors of $\Phi_{m_c}$ modulo $p^\varepsilon$ and reduce them modulo $p^j$; the CRT gives $\ell$ unramified degree-$d$ extension-ring slots.

Let $\tau_a(X)=X^a$ for $a\in(\mathbb Z/m_c\mathbb Z)^\times$. The trace pairing is nondegenerate modulo $p$, since the algebra is a product of finite fields; its Gram determinant is consequently a unit modulo every $p^j$. There is a trace-dual basis $b_0,\ldots,b_{n-1}$ such that
\[
\operatorname{Tr}(b_uX^v)=\delta_{uv},\qquad
\operatorname{Tr}(W)=\sum_a\tau_a(W),\qquad
z_u=\operatorname{Tr}(b_uZ).
\]
Every base-ring-linear map $\mathcal T:R_j\to R_j$ therefore has the explicit realization
\[
\mathcal T(Z)=\sum_a\kappa_a\tau_a(Z),\qquad
\kappa_a=\sum_{u=0}^{n-1}\mathcal T(X^u)\tau_a(b_u).
\]
Indeed, substituting the trace expansion gives $\sum_u\mathcal T(X^u)z_u$. Each term uses a public polynomial multiplier and an automorphism followed by key switching. No decryption or division by the extension degree is involved. In particular, when $p\mid d$, dividing the slot trace by $d$ is invalid. The lifted Frobenius is the automorphism $X\mapsto X^p$ acting on generators; it is not the map $W\mapsto W^p$ on arbitrary elements modulo $p^j$ when $j>1$.

For $0\leq u<d$, define $\mathcal T_u$ by giving its CRT residues: slot $h$ contains the scalar coefficient $z_{u\ell+h}$ and no nonconstant extension coordinates. An inverse map $\mathcal U_u$ reads each slot's scalar coordinate and writes it to coefficient position $u\ell+h$, filling other positions with zero. CRT matrix inversion constructs these maps, and the trace-dual formula compiles both into homomorphic operations. They satisfy
\[
\sum_{u=0}^{d-1}\mathcal U_u\bigl(\mathcal T_u(Z)\bigr)=Z.
\]
Use the packing maps modulo $p^\varepsilon$ and the compatible unpacking maps modulo $p^r$. Digit polynomials act on the scalar entries, not on arbitrary extension elements. If every slot of their numerator is divisible by $p$, the CRT implies that the entire polynomial lies in $pR_j$; thus the subsequent coefficientwise exact division is well defined. This supplies a general, unoptimized construction for the transforms used below~\cite{cryptoeprint:2022/1363}.

The familiar completely split matrices are the special case $d=1$:

The goal of the \textsf{CoeffToSlot} step is to homomorphically move polynomial $Z$'s coefficients $z_i$ to input vector slots. 


In Summary~\ref*{subsubsec:bfv-rotation-summary} (in \autoref{subsubsec:bfv-rotation-summary}), we learned that the encoding formula for converting a vector of input slots $\vec{v}$ into a vector of polynomial coefficients $\vec{m}$ is: $\vec{m} = n^{-1}\cdot\hathat{W} \cdot I_n^R \cdot \vec{v}$, where $\hathat{W}$ is a basis of the $n$-dimensional vector space crafted as follows: 

$\hathat{W} = \begin{bmatrix}
1 & 1 & \cdots & 1 & 1 & 1 & \cdots & 1\\
(\omega^{J(\frac{n}{2} - 1)}) & (\omega^{J(\frac{n}{2} - 2)}) & \cdots & (\omega^{J(0)}) & (\omega^{J_*(\frac{n}{2} - 1)}) & (\omega^{J_*(\frac{n}{2} - 2)}) & \cdots & (\omega^{J_*(0)})\\
(\omega^{J(\frac{n}{2} - 1)})^2 & (\omega^{J(\frac{n}{2} - 2)})^2 & \cdots & (\omega^{J(0)})^2 & (\omega^{J_*(\frac{n}{2} - 1)})^2 & (\omega^{J_*(\frac{n}{2} - 2)})^2 & \cdots & (\omega^{J_*(0)})^2 \\
\vdots & \vdots & \ddots & \vdots & \vdots & \vdots & \ddots & \vdots \\
(\omega^{J(\frac{n}{2} - 1)})^{n-1} & (\omega^{J(\frac{n}{2} - 2)})^{n-1} & \cdots & (\omega^{J(0)})^{n-1} & (\omega^{J_*(\frac{n}{2} - 1)})^{n-1} & (\omega^{J_*(\frac{n}{2} - 2)})^{n-1} & \cdots  & (\omega^{J_*(0)})^{n-1}
\end{bmatrix}$

$ $

\textcolor{red}{ $\rhd$ where the rotation helper function $J(h) = 5^h \bmod 2n$}

$ $

Therefore, given the input ciphertext $\textsf{ct} = \textsf{RLWE}_{S, \sigma}\bm(\Delta' Z\bm) \bmod q$, we can understand its input vector slots as storing some values such that multiplying $n^{-1}\cdot\hathat{W} \cdot I_n^R$ to each of them turns them into a coefficient $z_i$ of polynomial $Z$. This implies that if we \textit{homomorphically} multiply $n^{-1}\cdot\hathat{W} \cdot I_n^R$ to the input vector slots of $\textsf{RLWE}_{S, \sigma}\bm(\Delta' Z\bm)$, then the resulting ciphertext's $n$-slot input vector slots will contain the $n$ coefficients of $Z$, which is equivalent to moving the coefficients of $Z$ to the input vector slots. Therefore, the \textsf{CoeffToSlot} step is equivalent to homomorphically computing $n^{-1}\cdot\hathat{W} \cdot I_n^R \cdot \textsf{RLWE}_{S, \sigma}\bm(\Delta' Z\bm)$. We can homomorphically compute matrix-vector multiplication by using the technique explained in \autoref{subsec:bfv-matrix-multiplication}.

After the \textsf{CoeffToSlot} step, we can homomorphically eliminate the noise in the lower (base-$p$) $\varepsilon-r$ digits of each $z_i$ by homomorphically evaluating the noise-removing polynomial (to be explained in the next subsection).

After recovering the coefficients of $M$ in the plaintext slots, we need to move them back to their original coefficient positions. The ciphertext still contains bounded evaluation error. This step is called \textsf{SlotToCoeff}, which is an exact inverse procedure of \textsf{CoeffToSlot}. We also learned in Summary~\ref*{subsubsec:bfv-rotation-summary} (in \autoref{subsubsec:bfv-rotation-summary}) that the inverse matrix of $n^{-1}\cdot\hathat{W} \cdot I_n^R$ is $\hathat{W}^*$, where:

$\hathat{W}^* = \begin{bmatrix}
1 & (\omega^{J(0)}) & (\omega^{J(0)})^2 & \cdots & (\omega^{J(0)})^{n-1}\\
1 & (\omega^{J(1)}) & (\omega^{J(1)})^2 & \cdots & (\omega^{J(1)})^{n-1}\\
1 & (\omega^{J(2)}) & (\omega^{J(2)})^2 & \cdots & (\omega^{J(2)})^{n-1}\\
\vdots & \vdots & \vdots & \ddots & \vdots \\
1 & (\omega^{J(\frac{n}{2}-1)}) & (\omega^{J(\frac{n}{2}-1)})^2 & \cdots & (\omega^{J(\frac{n}{2}-1)})^{n-1}\\
1 & (\omega^{J_*(0)}) & (\omega^{J_*(0)})^2 & \cdots & (\omega^{J_*(0)})^{n-1}\\
1 & (\omega^{J_*(1)}) & (\omega^{J_*(1)})^2 & \cdots & (\omega^{J_*(1)})^{n-1}\\
1 & (\omega^{J_*(2)}) & (\omega^{J_*(2)})^2 & \cdots & (\omega^{J_*(2)})^{n-1}\\
\vdots & \vdots & \vdots & \ddots & \vdots \\
1 & (\omega^{J_*(\frac{n}{2}-1)}) & (\omega^{J_*(\frac{n}{2}-1)})^2 & \cdots & (\omega^{J_*(\frac{n}{2}-1)})^{n-1}\\
\end{bmatrix}$

Therefore, the \textsf{SlotToCoeff} step is equivalent to homomorphically multiplying $\hathat{W}^*$ to the output of the noise-eliminating polynomial evaluation.  

In the next subsection, we will learn how to design the core algorithm of BFV, the noise elimination polynomial, based on the digit extraction polynomial $G_\varepsilon(x)$. 


\subsubsection{Digit Extraction}
\label{subsubsec:bfv-bootstrapping-digit-extraction}


For odd $p$, use balanced base-$p$ digits $d\in\{-(p-1)/2,\ldots,(p-1)/2\}$. Write an input coefficient as $z=\sum_{j=0}^{\varepsilon-1}d_jp^j\pmod{p^\varepsilon}$. Deleting its lowest $v=\varepsilon-r$ balanced digits gives $\lceil z/p^v\rfloor\pmod{p^r}$, as required by decryption. Ordinary nonnegative digits would instead give flooring: for example $p=3,\varepsilon=3,m=1,e'=-1$ gives $z=8$, and $\lfloor8/9\rfloor=0\ne1$ despite $|e'|<9/2$. For $p=2$, use nonnegative digits and add $2^{v-1}$ before deleting $v$ bits, implementing $\lceil z/2^v\rfloor=\lfloor(z+2^{v-1})/2^v\rfloor$ away from ties. The $p=2$ formula is only a scalar rounding identity here. It does not extend the cited CRT bootstrapping route to this chapter's power-of-two ring; that route requires a cyclotomic index coprime to $p$, as explained in \autoref{subsubsec:bfv-bootstrapping-coefftoslot}.

For $j\geq1$, a lifting polynomial $F_j$ satisfies
\[
F_j(d+p^ju)\equiv d\pmod{p^{j+1}}
\]
for each allowed lowest digit $d$ and integer $u$. The index $j$ denotes a known precision in the algorithm; it need not be the first nonzero higher digit of a secret input. Let $L_\varepsilon=F_{\varepsilon-1}\circ\cdots\circ F_1$.

\para{1st Recursion:} $F_1(z)\equiv d_0\pmod{p^2}$.

\para{2nd Recursion:} $F_2\circ F_1(z)\equiv d_0\pmod{p^3}$.

\para{3rd Recursion:} $F_3\circ F_2\circ F_1(z)\equiv d_0\pmod{p^4}$.

\para{$\bm{(\varepsilon-1)}$-th Recursion:} $L_\varepsilon(z)\equiv d_0\pmod{p^\varepsilon}$.

The digit-removal operation is therefore
\[
G_\varepsilon(z)=\frac{z-L_\varepsilon(z)}p\pmod{p^{\varepsilon-1}}.
\]
This is exact integer division of a divisible representative, with the output modulus reduced by one power of $p$. It is not multiplication by an inverse of $p$ in $\mathbb Z_{p^\varepsilon}$, where such an inverse does not exist. Even if $p^{-1}\bmod q$ exists, multiplying ciphertext components by it is a different operation and can amplify noise. Thus $G_\varepsilon$ denotes a polynomial-evaluation-and-division procedure, not a single polynomial over $\mathbb Z_{p^\varepsilon}$. Repeating
\[
G_{r+1}\circ\cdots\circ G_{\varepsilon-1}\circ G_\varepsilon(z)
\]
deletes $\varepsilon-r$ digits. Each invocation of $G_j$ computes its own lowest digit using $F_1,\ldots,F_{j-1}$, starting at $F_1$ even when $r>1$.

\para{Scaling Factor Re-interpretation:} Suppose a homomorphic subtraction leaves a plaintext representative $u=pw\pmod{p^j}$ and ciphertext phase $\Delta_j u+E_j$, with $\Delta_j=\lfloor q/p^j\rfloor$. With a consistent divisible centered representative, the same ciphertext can be interpreted at plaintext modulus $p^{j-1}$ as
\[
\Delta_jpw+E_j=\Delta_{j-1}w+(E_j-\delta_jw),\qquad
\delta_j=\Delta_{j-1}-p\Delta_j\in\{0,\ldots,p-1\}.
\]
This avoids multiplying ciphertext components by $p^{-1}\bmod q$. The term $-\delta_jw$ is a change in the error measured against the floored scale. Against the exact scale it cancels:

\[
\frac{q}{p^j}(pw)+\mathcal E_j
=\frac{q}{p^{j-1}}w+\mathcal E_j.
\]
Thus exact division adds no natural error and increases the decoder tolerance. Both descriptions are equivalent: one must not add the floor-scale drift again to a bound already expressed in natural-error coordinates. Consistent divisible centered coefficient representatives make the identity exact, including when $p^j\nmid q$.

The floored scales are $\lfloor q/p^\varepsilon\rfloor,\lfloor q/p^{\varepsilon-1}\rfloor,\ldots,\lfloor q/p^r\rfloor$. A plaintext added to a ciphertext must use the current scale, while ciphertext-to-plaintext multiplication uses an unscaled multiplier as in \autoref{subsec:bfv-mult-plain}. Polynomial evaluation changes the natural error; exact division preserves it while changing the plaintext modulus. The explicit recurrence in \autoref{subsubsec:bfv-bootstrapping-summary} checks every intermediate decoder condition.

The remaining local algebraic task is to construct the lifting polynomials.

\para{Designing \boldmath$F_{\varepsilon'}(z_i)$:} We give a precision-aware interpolation construction. Let $j=\varepsilon'\geq1$ and let $D$ be the chosen set of $p$ consecutive digit representatives. For every $d\in D$, Fermat's theorem makes $(d^p-d)/p$ an integer. Interpolate a polynomial $H_j$ of degree at most $p-1$ over $\mathbb Z_{p^j}$ with
\[
H_j(d)\equiv(d^p-d)/p\pmod{p^j}\quad(d\in D).
\]
Every difference between distinct digit nodes is a unit modulo $p^j$, so Lagrange interpolation is valid. Choose integer lifts of the coefficients and define
\[
F_j(X)=X^p-pH_j(X).
\]
\begin{proof}
For $z=d+p^ju$, the binomial theorem gives $z^p\equiv d^p\pmod{p^{j+1}}$. Integer-coefficient polynomial evaluation also gives $H_j(z)\equiv H_j(d)\pmod{p^j}$. Hence
\[
F_j(z)\equiv d^p-pH_j(d)\equiv d\pmod{p^{j+1}}.
\]
\end{proof}
This proves the lifting property used in each recursion. The degree is at most $p$, rather than $p-1$. Interpolating only the separate digit functions over $\mathbb F_p$ and reusing arbitrary integer coefficient lifts is insufficient at higher powers: carries affect the higher digits. The coefficients must be constructed at the required precision, as above. This local existence proof does not by itself bound the depth or noise of an optimized homomorphic implementation.

\subsubsection{Summary}
\label{subsubsec:bfv-bootstrapping-summary}

The following summarizes the generalized plaintext transformation. The recurrence after the box makes its correctness conditions explicit: it accounts for the low-modulus rounding, the general linear transforms, all homomorphic gates and the final decoder. Security and optimization are separate questions.

\begin{tcolorbox}[title={\textbf{\tboxlabel{\ref*{subsubsec:bfv-bootstrapping-summary}} BFV Bootstrapping}}]
Let $t=p^r$, $\varepsilon>r$, $P=p^\varepsilon$, $\Delta=\lfloor q/p^r\rfloor$ and $\Delta'=\lfloor q/P\rfloor$. The complete CRT procedure assumes the compatible ring conditions in \autoref{subsubsec:bfv-bootstrapping-coefftoslot}; in this chapter's power-of-two ring, $p$ is odd.
\begin{enumerate}
\item \textbf{Modulus Switch:} Compute $A'=[\lceil(P/q)A\rfloor]_P$ and $B'=[\lceil(P/q)B\rfloor]_P$. Their phase satisfies $A'S+B'\equiv p^{\varepsilon-r}M+E'\pmod P$, with the complete $E'$ from \autoref{subsubsec:bfv-bootstrapping-modulus-switch}. Require $\|E'\|_\infty<p^{\varepsilon-r}/2$.
\item \textbf{Homomorphic Decryption:} With an encryption of $\Delta'S$ at modulus $q$, compute $A'\cdot\textsf{RLWE}_{S,\sigma}(\Delta'S)+(0,\Delta'B')$. Track both the key-noise term and the plaintext-reduction remainder as in \autoref{subsubsec:bfv-bootstrapping-homomorphic-decryption}.
\item \textbf{\textsf{CoeffToSlot}:} Homomorphically place the coefficients of $Z=[p^{\varepsilon-r}M+E']_P$ into slots. In the scalar-split case the slot transformation is $n^{-1}\hathat W I_n^R$ (\autoref{subsubsec:bfv-rotation-summary}); for extension slots, use the $d$ maps $\mathcal T_u$ constructed in \autoref{subsubsec:bfv-bootstrapping-coefftoslot}.
\item \textbf{Digit Extraction:} For each current precision $j=\varepsilon,\varepsilon-1,\ldots,r+1$, extract its lowest digit with $L_j=F_{j-1}\circ\cdots\circ F_1$ and implement the exact quotient $G_j(z)=(z-L_j(z))/p\pmod{p^{j-1}}$. Use balanced digits for odd $p$, or the stated rounding offset for $p=2$. The final slot is $m_i\pmod{p^r}$, so any unreduced representative differs by $k_i'p^r$, not by $k_i'p$ when $r>1$.
\para{\underline{Scaling Factor Re-interpretation}:} At each division, reinterpret the scale from $\lfloor q/p^j\rfloor$ to $\lfloor q/p^{j-1}\rfloor$ and include the drift $-\delta_jw$ derived above. The final scale is $\lfloor q/p^r\rfloor=\Delta$. No modular inverse of $p$ is assumed.
\item \textbf{\textsf{SlotToCoeff}:} Apply the inverse linear transform $\hathat W^*$ in the scalar-split case, or sum the $\mathcal U_u$ outputs in the extension-slot case, recovering an encryption with phase $\Delta M+E^{\langle b\rangle}\pmod q$. The complete output error includes all evaluation, scale-drift and plaintext-reduction terms and must satisfy the BFV decoder bound.
\end{enumerate}
\end{tcolorbox}

\textbf{An explicit conditional correctness certificate.} Use the natural error $\mathcal E_T$ defined in \autoref{subsubsec:bfv-bootstrapping-high-level}. Let $C_R$ satisfy $\|UV\|_\infty\leq C_R\|U\|_\infty\|V\|_\infty$, and let $h_S$ bound the infinity operator norm of multiplication by the secret. In the negacyclic ring, one may take $C_R=n$ and $h_S=\|S\|_1$. For a public polynomial $C$ or automorphism $\tau_a$, let $\lambda_C$ or $\lambda_a$ be its coefficient infinity operator norm. These are computable matrix norms, also for general cyclotomic rings.

Use exact integer gadgets with $L$ digits bounded by $D_g$, and raw evaluation-key row phases $g_iS_{\rm old}+E_i^{\rm key}$ under the output secret, with $\|E_i^{\rm key}\|_\infty\leq B_{\rm key}$. A key switch adds at most
\[
B_{\rm ks}=C_RLD_gB_{\rm key}
\]
to the natural-error bound. No independence assumption is needed. If input bounds are $B_1,B_2$, the following valid bounds can be propagated, rounding them upward if desired:
\[
\begin{aligned}
B_{\rm add/sub}&=B_1+B_2,\qquad B_{\rm plain\,mult}=\lambda_CB_1,\\
B_{\rm plain\,add}&=B_1+(r_T/T)\|C\|_\infty,\\
B_{\rm auto}&=\lambda_aB_1+B_{\rm ks},\qquad B_{\rm exact\,div}=B_1,\\
\kappa_i&=(h_S+2)/2+B_i/q.
\end{aligned}
\]
The plaintext addition uses $\lfloor q/T\rfloor C$; exact division requires a plaintext divisible by $p$ and changes $T$ to $T/p$. Standard FV multiplication gives
\[
\begin{aligned}
B_{\rm mult}={}&C_R\left[\frac T2(B_1+B_2)
+T(\kappa_1B_2+\kappa_2B_1)+\frac TqB_1B_2\right]\\
&+\frac{1+h_S+h_S^2}{2}+B_{\rm ks}.
\end{aligned}
\]
To verify this bound, write centered component lifts as $AS+B=(q/T)M+\mathcal E+qK$. Their norms give $\|K\|_\infty\leq\kappa$. Expanding the integer tensor product and multiplying by $T/q$ leaves the natural-error terms
\[
\begin{aligned}
\mathcal E_{\rm mult}={}&M_1\mathcal E_2+M_2\mathcal E_1
+T(\mathcal E_1K_2+\mathcal E_2K_1)+\frac Tq\mathcal E_1\mathcal E_2\\
&+\rho_0+\rho_1S+\rho_2S^2+E_{\rm ks}\pmod q.
\end{aligned}
\]
All omitted carry terms are integer multiples of $q$. Each tensor-component rounding error $\rho_i$ has norm at most $1/2$, giving the displayed rounding bound. Reducing the message modulo $T$ also changes the exact scaled message only by multiples of $q$. This proves the recurrence, rather than assuming that multiplication is correct.

For the first modulus switch, an input natural-error bound $B_{\rm in}$ is sufficient if
\[
\frac PqB_{\rm in}+\frac{1+h_{\rm in}}2
<\frac{p^{\varepsilon-r}}2.
\]
Here $h_{\rm in}$ bounds multiplication by the input secret. The old error is now part of the integer plaintext $Z$; it is not the error of the new evaluation ciphertext. If the encryption of the input secret has natural-error bound $B_s$, the homomorphic decryption starts the evaluation ledger with
\[
B_Z\leq\lambda_{A'}B_s+(r_P/P)\|B'\|_\infty.
\]
For a floor-scaled bootstrapping-key encryption with integer error bounded by $B_{\rm key}$, one may use $B_s=B_{\rm key}+(r_P/P)\|S_{\rm in}\|_\infty$. The evaluation modulus may be chosen independently of the input modulus; the displayed formulas use the same $q$ for simplicity.

Require the bound $B_C<q/(2T_C)$ at every ciphertext node of the explicit circuit. The modulus-switch inequality selects the correct rounded message; the trace-dual maps place its coefficients in scalar slots; the lifting polynomials and exact divisions recover those coefficients modulo $p^r$; the inverse maps reconstruct $M$. Induction with the above recurrence proves that the final ciphertext decrypts to $M$. In floor-scale notation its output error is $E^{\langle b\rangle}=\mathcal E_{p^r}+(r_{p^r}/p^r)M$. The new bound need not be smaller than every possible input error; it must fit the intended subsequent computation.

The exact public-evaluation fixtures in \texttt{audit/check\_bootstrap.py} cover $\Phi_4,p=5$ with scalar splitting, $\Phi_8,p=3$ with degree-two extension slots at $r=1$ and $r=2$, and the coprime binary alternative $\Phi_3,p=2$. Their ciphertext moduli are not divisible by the plaintext moduli, and their input and evaluation-key errors are nonzero. A separate checker audits every intermediate ciphertext after public evaluation; no secret or decryption callback enters the evaluator. These tiny-ring fixtures verify the implementation and algebra, not cryptographic security or the optimized thin-bootstrapping algorithms of the reference.

\para{The purpose of Homomorphic Decryption:} It converts a low-modulus ciphertext's secret-dependent phase into an encrypted plaintext at a fresh evaluation modulus and suitable scale. This makes the remaining rounding circuit evaluable without revealing the secret or message. Scale reinterpretation is one component of the digit-removal implementation, not by itself a complete explanation or correctness proof of bootstrapping.

